% ============================================================
% Capítulo 3 - Diseño y desarrollo de la solución
% Investigación en ingeniería
%
% Paquetes necesarios en el preámbulo del proyecto:
%   \usepackage[utf8]{inputenc}   \usepackage[spanish]{babel}
%   \usepackage{booktabs, tabularx, array, graphicx, amsmath, enumitem, xcolor}
%
% Cada sección vive en su propio archivo (seccion_3_X.tex). Copia los 6 archivos
% a tu proyecto y usa % ============================================================
% Capítulo 3 - Diseño y desarrollo de la solución
% Investigación en ingeniería
%
% Paquetes necesarios en el preámbulo del proyecto:
%   \usepackage[utf8]{inputenc}   \usepackage[spanish]{babel}
%   \usepackage{booktabs, tabularx, array, graphicx, amsmath, enumitem, xcolor}
%
% Cada sección vive en su propio archivo (seccion_3_X.tex). Copia los 6 archivos
% a tu proyecto y usa % ============================================================
% Capítulo 3 - Diseño y desarrollo de la solución
% Investigación en ingeniería
%
% Paquetes necesarios en el preámbulo del proyecto:
%   \usepackage[utf8]{inputenc}   \usepackage[spanish]{babel}
%   \usepackage{booktabs, tabularx, array, graphicx, amsmath, enumitem, xcolor}
%
% Cada sección vive en su propio archivo (seccion_3_X.tex). Copia los 6 archivos
% a tu proyecto y usa % ============================================================
% Capítulo 3 - Diseño y desarrollo de la solución
% Investigación en ingeniería
%
% Paquetes necesarios en el preámbulo del proyecto:
%   \usepackage[utf8]{inputenc}   \usepackage[spanish]{babel}
%   \usepackage{booktabs, tabularx, array, graphicx, amsmath, enumitem, xcolor}
%
% Cada sección vive en su propio archivo (seccion_3_X.tex). Copia los 6 archivos
% a tu proyecto y usa \input{cap3} desde el documento principal.
% ============================================================

% Marca visible para lo que el equipo debe completar o confirmar.
\newcommand{\pendiente}[1]{\textbf{\textcolor{red}{[POR DEFINIR: #1]}}}

\chapter{Diseño y desarrollo de la solución}
\label{cap:diseno}

En este capítulo se describe el proceso de diseño y desarrollo de un sistema de autenticación continua basado en la señal electrocardiográfica (ECG) adquirida con una banda Polar H10. Se presentan los requerimientos de la solución, las alternativas tecnológicas consideradas, el diseño detallado del sistema, la construcción del prototipo y las características de la solución finalmente desarrollada.

El proyecto se llevó a cabo mediante una metodología de prototipado iterativo. Primero se construyó un extractor de embeddings entrenado con una base de datos pública, después se incorporó la adquisición de señal con el sensor y el protocolo de recolección con voluntarios, y finalmente se integraron los módulos de decisión continua. Cada iteración permitió identificar problemas y mejorar el sistema. Las tablas de resultados, el desempeño medido y la comprobación de la hipótesis se presentan en el Capítulo~\ref{cap:resultados}.

\input{seccion_3_1}
\input{seccion_3_2}
\input{seccion_3_3}
\input{seccion_3_4}
\input{seccion_3_5}

\section*{Cierre del capítulo}

En este capítulo se definieron los requerimientos funcionales, de desempeño y no funcionales del sistema; se generaron y compararon alternativas mediante técnicas de creatividad y una matriz de decisión; y se seleccionó una arquitectura basada en una red neuronal convolucional unidimensional con bloques residuales que produce embeddings de identidad. También se describieron el diseño de la red, el preprocesamiento, la estrategia de entrenamiento y la partición de datos por persona, así como la implementación modular del prototipo y su protocolo de recolección. En el siguiente capítulo se presentan los resultados de las pruebas experimentales, el análisis de los indicadores de desempeño y la comprobación del cumplimiento de los objetivos y de la hipótesis.
 desde el documento principal.
% ============================================================

% Marca visible para lo que el equipo debe completar o confirmar.
\newcommand{\pendiente}[1]{\textbf{\textcolor{red}{[POR DEFINIR: #1]}}}

\chapter{Diseño y desarrollo de la solución}
\label{cap:diseno}

En este capítulo se describe el proceso de diseño y desarrollo de un sistema de autenticación continua basado en la señal electrocardiográfica (ECG) adquirida con una banda Polar H10. Se presentan los requerimientos de la solución, las alternativas tecnológicas consideradas, el diseño detallado del sistema, la construcción del prototipo y las características de la solución finalmente desarrollada.

El proyecto se llevó a cabo mediante una metodología de prototipado iterativo. Primero se construyó un extractor de embeddings entrenado con una base de datos pública, después se incorporó la adquisición de señal con el sensor y el protocolo de recolección con voluntarios, y finalmente se integraron los módulos de decisión continua. Cada iteración permitió identificar problemas y mejorar el sistema. Las tablas de resultados, el desempeño medido y la comprobación de la hipótesis se presentan en el Capítulo~\ref{cap:resultados}.

\section{Especificación de requerimientos}
\label{sec:requerimientos}

Esta sección establece qué debe hacer el sistema y bajo qué condiciones debe funcionar. Los requerimientos sirven como referencia para diseñar la solución y para comprobar, en el Capítulo~\ref{cap:resultados}, si el prototipo cumple con lo solicitado.

\subsection{Descripción general del sistema}

El sistema simula un examen en computadora. Mientras la persona lo contesta, una banda Polar H10 colocada en el pecho registra su ECG. Una red neuronal convierte ventanas sucesivas de la señal en vectores de características (\emph{embeddings}) y los compara con la plantilla creada en el registro inicial del usuario. De esta forma el sistema verifica de manera continua que quien contesta el examen es la misma persona que se registró, y no solo en el momento de iniciar sesión.

\subsection{Requerimientos funcionales}

\begin{itemize}
  \item Adquirir la señal ECG del Polar H10 por Bluetooth Low Energy a 130~Hz.
  \item Verificar la calidad de la señal antes y durante la grabación (señal plana, ruido excesivo, picos anómalos, latidos insuficientes).
  \item Preprocesar la señal: filtro pasa-banda de 0.5 a 40~Hz, segmentación en ventanas de 5~s con 50\,\% de traslape y normalización z-score por ventana.
  \item Generar un embedding L2-normalizado de 128 dimensiones por ventana.
  \item Crear la plantilla de un usuario durante el registro (sesión de enrolamiento) a partir del promedio de sus embeddings.
  \item Calcular la similitud coseno entre cada ventana nueva y la plantilla, y decidir si la identidad se acepta o se rechaza.
  \item Aplicar una regla de decisión continua: los rechazos aislados no se muestran al usuario; tras 3 rechazos consecutivos se le pide una nueva muestra (por ejemplo, sentarse derecho y permanecer quieto unos segundos) y, si esta tampoco es aceptada, se bloquea el examen.
  \item Entrenar y evaluar el modelo con personas que no participaron en el entrenamiento.
  \item Registrar los datos de cada sesión con identificadores anónimos, para poder repetir los análisis.
\end{itemize}

\subsection{Requerimientos no funcionales y restricciones}

\begin{itemize}
  \item El entrenamiento debe ejecutarse en una computadora sin GPU (16~GB de RAM, CPU de 6 núcleos), en un tiempo del orden de minutos a una hora.
  \item Los experimentos deben ser reproducibles: semilla fija, un único archivo de configuración y un comando por etapa de entrenamiento.
  \item La red neuronal debe permanecer separada de la base de plantillas y de la lógica de decisión, de modo que cada componente pueda sustituirse sin modificar los demás.
  \item Los datos fisiológicos deben tratarse de forma anónima y no deben publicarse en el repositorio.
  \item El modelo debe poder explicarse por completo en cada una de sus partes por los integrantes del equipo.
  \item Restricciones del proyecto: una muestra de 20 voluntarios, un solo canal de ECG, un solo sensor disponible y un plazo de desarrollo que concluye en noviembre de 2026.
\end{itemize}

\subsection{Tabla de requerimientos}

\begin{table}[htbp]
\centering
\caption{Especificación de requerimientos y criterios de aceptación.}
\label{tab:requerimientos}
\small
\begin{tabularx}{\textwidth}{l X l X}
\toprule
\textbf{Código} & \textbf{Requerimiento} & \textbf{Tipo} & \textbf{Criterio de aceptación} \\
\midrule
RF-01 & Adquirir ECG del Polar H10 a 130~Hz & Funcional & Se guardan sesiones completas sin pérdida de conexión \\
RF-02 & Evaluar la calidad de la señal & Funcional & Una sesión con señal defectuosa se detecta y se repite \\
RF-03 & Generar un embedding por ventana de 5~s & Funcional & Salida de dimensión 128 con norma L2 igual a 1 \\
RF-04 & Crear la plantilla de cada usuario & Funcional & Plantilla construida solo con la sesión de enrolamiento \\
RF-05 & Decidir identidad por similitud coseno & Funcional & Umbral fijado solo con datos de validación \\
RF-06 & Aplicar la regla de decisión continua & Funcional & Rechazos aislados ocultos; con 3 consecutivos se pide nueva muestra y, si falla, se bloquea \\
RD-01 & Distinguir personas nunca vistas por la red & Desempeño & EER medido en personas \emph{hold-out}, claramente por debajo del azar (50\,\%) y con intervalo de confianza que no lo incluya \\
RD-02 & Mantener el desempeño entre días & Desempeño & EER entre sesiones de días distintos reportado con intervalo de confianza \\
RD-03 & Detectar a un impostor durante el examen & Desempeño & Tras 3 rechazos consecutivos (unos 10~s) se solicita una nueva muestra; si esta también falla, se bloquea el examen \\
RD-04 & No interrumpir innecesariamente al usuario legítimo & Desempeño & Los rechazos aislados no se notifican; la tasa de bloqueos injustificados se reporta en el Capítulo~\ref{cap:resultados} \\
RNF-01 & Entrenar sin GPU & No funcional & Entrenamiento completo en CPU en menos de una hora \\
RNF-02 & Reproducibilidad & No funcional & Misma semilla y configuración producen el mismo modelo \\
RNF-03 & Separación entre red, plantillas y decisión & No funcional & Cada módulo se prueba de forma independiente \\
RNF-04 & Anonimato de los participantes & No funcional & Identificadores fijos tipo P01 y datos excluidos del repositorio \\
\bottomrule
\end{tabularx}
\end{table}

No se fija una meta numérica rígida de EER: con 20 participantes y un solo canal, el objetivo es demostrar la viabilidad del enfoque. El criterio es que el EER quede claramente por debajo del azar y se compare con los valores reportados en el estado del arte del Capítulo~\ref{cap:marco}. Por ello, los resultados del proyecto se consideran una prueba de concepto.

\section{Generación, evaluación y selección de alternativas}
\label{sec:alternativas}

Esta sección muestra cómo se generaron distintas propuestas para resolver el problema y por qué se eligió una de ellas. Las técnicas de creatividad no son las alternativas finales: son procedimientos para generar ideas que después se organizan, comparan y seleccionan.

% NOTA PARA EL EQUIPO: las ideas, tablas y calificaciones de esta sección son una
% PROPUESTA de partida. Deben discutirse y ajustarse según lo que el equipo realmente
% consideró, y sustentarse con literatura o pruebas preliminares. Conserven evidencias
% (fotos de la lluvia de ideas, mapa mental) para anexarlas.

\subsection{Técnicas para la generación de ideas}

\subsubsection{Lluvia de ideas}

Pregunta generadora: ¿qué sensor, qué representación de la señal y qué modelo podrían utilizarse para verificar de forma continua la identidad de una persona a partir de su ECG? Las ideas obtenidas fueron: banda Polar H10, un sensor de ECG de diseño propio, otros equipos comerciales de mayor costo (por ejemplo, Movesense o Zephyr BioHarness), señal cruda o latidos segmentados, rasgos manuales (intervalos y amplitudes de las ondas P, QRS y T), máquina de vectores de soporte (SVM), red neuronal convolucional 1D como clasificador, red 1D con embeddings, pérdida de tripleta, preentrenamiento con una base de datos pública y suavizado de las decisiones en el tiempo. \pendiente{confirmar con el equipo qué equipo comercial más caro se consideró}

\subsubsection{Mapa mental}

Se colocó en el centro el concepto <<autenticación continua con ECG>> y se crearon ramas para adquisición, preprocesamiento, representación de la señal, modelo, entrenamiento, decisión continua, interfaz y validación. El mapa permitió verificar que las propuestas contemplaran el sistema completo y no solamente la red neuronal. \pendiente{insertar figura del mapa mental}

\subsubsection{SCAMPER}

\begin{table}[htbp]
\centering
\caption{Aplicación de SCAMPER al sistema de autenticación continua.}
\label{tab:scamper}
\small
\begin{tabularx}{\textwidth}{l X X}
\toprule
\textbf{Acción} & \textbf{Pregunta aplicada} & \textbf{Ejemplo de idea} \\
\midrule
Sustituir & ¿Qué componente puede reemplazarse? & Sustituir rasgos manuales por rasgos aprendidos por una red \\
Combinar & ¿Qué elementos pueden trabajar juntos? & Combinar preentrenamiento público con ajuste fino con datos del Polar \\
Adaptar & ¿Qué solución externa puede aprovecharse? & Adaptar el aprendizaje de métricas usado en reconocimiento facial \\
Modificar & ¿Qué característica puede cambiarse? & Cambiar la longitud de la ventana de 5~s \\
Otros usos & ¿Puede un elemento cumplir otra función? & Usar la penúltima capa de un clasificador como huella de identidad \\
Eliminar & ¿Qué elemento no es indispensable? & Eliminar la detección explícita de picos R y trabajar sobre la señal filtrada \\
Reordenar & ¿Qué ocurre al cambiar el orden? & Suavizar decisiones en el tiempo en lugar de decidir por ventana aislada \\
\bottomrule
\end{tabularx}
\end{table}

\subsubsection{Combinación de elementos y matriz morfológica}

La matriz morfológica combina opciones de distintos subsistemas para construir alternativas completas y técnicamente coherentes.

\begin{table}[htbp]
\centering
\caption{Matriz morfológica del sistema.}
\label{tab:morfologica}
\small
\begin{tabularx}{\textwidth}{l X X X}
\toprule
\textbf{Función} & \textbf{Opción 1} & \textbf{Opción 2} & \textbf{Opción 3} \\
\midrule
Adquisición & Polar H10 & Sensor de ECG de diseño propio & Equipo comercial de mayor costo \\
Entrada del modelo & Rasgos manuales & Ventana cruda filtrada & Latidos segmentados \\
Modelo & SVM & CNN 1D como clasificador & CNN 1D con embeddings \\
Pérdida & Entropía cruzada & Tripleta & ArcFace \\
Datos de entrenamiento & Solo Polar & Preentrenado y ajustado & Solo base pública \\
Decisión & Por ventana & Promedio móvil & Acumulación de evidencia \\
\bottomrule
\end{tabularx}
\end{table}

\subsection{Alternativas de solución obtenidas}

\begin{table}[htbp]
\centering
\caption{Alternativas conceptuales.}
\label{tab:alternativas}
\small
\begin{tabularx}{\textwidth}{l X X X}
\toprule
\textbf{Alternativa} & \textbf{Configuración} & \textbf{Ventaja principal} & \textbf{Limitación principal} \\
\midrule
A: clásica & Polar H10, rasgos manuales, SVM, decisión por ventana & Pocos datos necesarios y fácil de explicar & Depende de detectar con precisión las ondas; menor capacidad de generalizar a personas nuevas \\
B: embeddings con preentrenamiento & Polar H10, ventana filtrada, CNN 1D residual con embeddings, preentrenamiento público y ajuste fino, entropía cruzada & Aprende rasgos automáticamente; permite registrar personas nuevas sin reentrenar; el preentrenamiento compensa pocos voluntarios & Diferencia entre la base pública y el Polar; requiere más tiempo de entrenamiento \\
C: aprendizaje de métricas avanzado & Igual que B pero con pérdida de tripleta o ArcFace y detección de picos R & Mayor separación potencial entre identidades & Mayor complejidad y sensibilidad a hiperparámetros con pocos datos \\
\bottomrule
\end{tabularx}
\end{table}

\subsection{Criterios de evaluación}

\begin{table}[htbp]
\centering
\caption{Criterios de evaluación y ponderación.}
\label{tab:criterios}
\small
\begin{tabularx}{\textwidth}{l l X}
\toprule
\textbf{Criterio} & \textbf{Ponderación} & \textbf{Justificación} \\
\midrule
Capacidad de generalizar a personas nuevas & 30\,\% & El sistema debe funcionar con usuarios que la red nunca vio \\
Viabilidad con pocos datos & 20\,\% & Solo se cuenta con 20 voluntarios, una cantidad pequeña para entrenar redes profundas \\
Costo computacional & 15\,\% & Solo se dispone de CPU \\
Facilidad de implementación y explicación & 20\,\% & El equipo debe dominar cada componente y cumplir el plazo \\
Aprovechamiento del sensor disponible & 15\,\% & El Polar H10 ya está en poder del equipo \\
\bottomrule
\end{tabularx}
\end{table}

Cada alternativa se califica del 1 al 5, donde 1 es un desempeño muy deficiente y 5 uno excelente. El resultado ponderado se obtiene multiplicando la calificación por la ponderación; la suma permite comparar las alternativas.

\begin{table}[htbp]
\centering
\caption{Matriz de decisión (calificaciones preliminares, por justificar).}
\label{tab:decision}
\small
\begin{tabular}{l c c c c}
\toprule
\textbf{Criterio} & \textbf{Pond.} & \textbf{A} & \textbf{B} & \textbf{C} \\
\midrule
Generalización a personas nuevas & 0.30 & 2 & 4 & 5 \\
Viabilidad con pocos datos & 0.20 & 4 & 4 & 2 \\
Costo computacional & 0.15 & 5 & 4 & 3 \\
Facilidad de implementación y explicación & 0.20 & 4 & 4 & 2 \\
Aprovechamiento del sensor & 0.15 & 4 & 4 & 4 \\
\midrule
\textbf{Total ponderado} & 1.00 & 3.55 & 4.00 & 3.35 \\
\bottomrule
\end{tabular}
\end{table}

\pendiente{justificar cada calificación con literatura, pruebas preliminares o cálculos; las cifras mostradas son un borrador}

\subsection{Selección y justificación}

Después de aplicar las técnicas de creatividad se obtuvieron tres alternativas de solución, evaluadas según su capacidad de generalizar a personas nuevas, su viabilidad con pocos datos, su costo computacional, su facilidad de implementación y el aprovechamiento del sensor disponible. La alternativa B obtuvo la mayor puntuación porque ofrece un equilibrio adecuado entre capacidad de generalización y viabilidad: el uso de embeddings permite registrar usuarios nuevos solo calculando su plantilla, y el preentrenamiento con una base de datos pública mitiga la escasez de voluntarios. La alternativa C se consideró una mejora futura en caso de que el error obtenido con B no fuera suficiente, por lo que la pérdida de tripleta y ArcFace se dejaron deliberadamente fuera de esta versión. En cuanto al sensor, se consideró construir uno propio, pero se encontró que el Polar H10 es más preciso, por lo que se descartó. Los equipos comerciales de mayor costo no ofrecían una ventaja que justificara su precio, y el Polar H10 cuenta con estudios que lo validan frente a un ECG de referencia. Por lo anterior, se seleccionó la alternativa B, con el Polar H10, como base para el diseño detallado.

\section{Modelado, cálculos, simulación y diseño detallado}
\label{sec:diseno}

Esta sección explica técnicamente cómo se diseñó la alternativa seleccionada antes de construirla, de modo que otra persona pueda comprender la arquitectura del sistema y reproducir el diseño.

\subsection{Arquitectura general}

El sistema se organiza como una cadena de bloques independientes:
\begin{center}
\emph{Polar H10} $\rightarrow$ adquisición BLE $\rightarrow$ control de calidad $\rightarrow$ filtrado y ventaneo $\rightarrow$ red de embeddings $\rightarrow$ similitud coseno con la plantilla $\rightarrow$ decisión
\end{center}
\pendiente{insertar diagrama de bloques}

La red (los pesos) se mantiene separada de la base de plantillas y de la lógica de decisión. Así, registrar a un usuario nuevo solo requiere calcular su plantilla, sin reentrenar la red.

\subsection{Adquisición y preprocesamiento de la señal}

La banda Polar H10 entrega el ECG a $f_s = 130$~Hz. La señal se procesa de la siguiente forma:
\begin{enumerate}
  \item \textbf{Filtro pasa-banda} de 0.5 a 40~Hz, que elimina la deriva de la línea base y atenúa el ruido muscular de alta frecuencia. La frecuencia de Nyquist es de 65~Hz, por lo que la banda es válida.
  \item \textbf{Ventaneo} en segmentos de 5~s, es decir, $5 \times 130 = 650$ muestras, con 50\,\% de traslape.
  \item \textbf{Normalización z-score por ventana}, para eliminar diferencias de amplitud debidas al contacto de los electrodos:
  \begin{equation}
    \tilde{x}_i = \frac{x_i - \mu}{\sigma}
  \end{equation}
  donde $\mu$ y $\sigma$ son la media y la desviación estándar de la ventana.
\end{enumerate}
Cada ventana es un tensor de forma $(1, 650)$. El sistema no detecta explícitamente los picos R para el modelo: la red extrae por sí misma las características morfológicas y temporales.

\subsection{Red neuronal de embeddings}

\subsubsection{Tipo de red}

Se utiliza una red neuronal convolucional unidimensional (1D-CNN) con bloques residuales, usada como \textbf{extractor de embeddings} (aprendizaje de métricas) y no como clasificador final. Su salida es un vector de 128 dimensiones con norma L2 igual a 1, de modo que la similitud coseno entre dos ventanas indica si provienen de la misma persona.

\subsubsection{Arquitectura capa por capa}

\begin{table}[htbp]
\centering
\caption{Arquitectura de la red para una ventana de entrada $(1, 650)$.}
\label{tab:arquitectura}
\small
\begin{tabularx}{\textwidth}{c l X c c c}
\toprule
\textbf{\#} & \textbf{Bloque} & \textbf{Operación} & \textbf{Kernel / stride} & \textbf{Canales} & \textbf{Salida} \\
\midrule
0 & Entrada & Ventana z-score & -- & 1 & $1 \times 650$ \\
1 & Stem & Conv1d + BatchNorm + ReLU & 7 / 1 & $1 \rightarrow 16$ & $16 \times 650$ \\
2 & Etapa 1 & Bloque residual & 5 / 2 & $16 \rightarrow 32$ & $32 \times 325$ \\
3 & Etapa 2 & Bloque residual & 5 / 2 & $32 \rightarrow 64$ & $64 \times 163$ \\
4 & Etapa 3 & Bloque residual & 5 / 2 & $64 \rightarrow 128$ & $128 \times 82$ \\
5 & Etapa 4 & Bloque residual & 5 / 2 & $128 \rightarrow 128$ & $128 \times 41$ \\
6 & Pooling global & Promedio sobre el tiempo & -- & 128 & 128 \\
7 & Embedding & Capa densa & -- & $128 \rightarrow 128$ & 128 \\
8 & Normalización & L2 & -- & 128 & 128 (norma 1) \\
9 & Cabeza & Capa densa sobre el embedding sin normalizar (solo entrenamiento) & -- & $128 \rightarrow N$ & $N$ \\
\bottomrule
\end{tabularx}
\end{table}

Todas las convoluciones usan \texttt{bias=False}, pues la normalización por lotes (BatchNorm) ya aporta el desplazamiento. La red tiene \textbf{371{,}024 parámetros} en el codificador; la cabeza agrega $128N + N$ parámetros, con $N$ igual al número de personas de entrenamiento, y se descarta al terminar.

\subsubsection{Bloque residual}

Cada etapa usa el mismo bloque. La rama principal aplica una convolución con $k=5$ y \emph{stride} 2, BatchNorm y ReLU, y después una segunda convolución con $k=5$ y \emph{stride} 1 seguida de BatchNorm. El atajo (\emph{skip connection}) es una convolución de $1\times1$ con \emph{stride} 2 seguida de BatchNorm, para que las dimensiones coincidan en la suma. La ReLU se aplica después de sumar ambas ramas:
\begin{equation}
  \mathbf{y} = \mathrm{ReLU}\big(F(\mathbf{x}) + S(\mathbf{x})\big)
\end{equation}
donde $F$ es la rama principal y $S$ el atajo. El \emph{stride} 2 reduce a la mitad la longitud temporal en cada etapa ($650 \rightarrow 325 \rightarrow 163 \rightarrow 82 \rightarrow 41$); con ello el campo receptivo de la última etapa abarca varios latidos.

\subsubsection{Embedding y similitud}

El pooling global promedio hace que la huella sea independiente de en qué punto de la ventana cae cada latido. La salida se normaliza:
\begin{equation}
  \mathbf{e} = \frac{\mathbf{z}}{\lVert \mathbf{z} \rVert_2}
\end{equation}
y la similitud entre una ventana con embedding $\mathbf{e}$ y una plantilla $\mathbf{t}$ es el coseno del ángulo entre ambos, que al estar normalizados equivale al producto punto:
\begin{equation}
  s(\mathbf{e}, \mathbf{t}) = \mathbf{e} \cdot \mathbf{t}
\end{equation}
La plantilla de una persona es el promedio de los embeddings de su sesión de enrolamiento. La identidad se acepta si $s$ supera un umbral $\tau$.

\subsection{Estrategia de entrenamiento}

\begin{enumerate}
  \item \textbf{Preentrenamiento} con la base pública ECG-ID, remuestreada a 130~Hz: la red se entrena como clasificador de identidades con entropía cruzada y \emph{label smoothing} de 0.1. Su propósito es aprender rasgos generales del ECG.
  \item \textbf{Ajuste fino} con los voluntarios del Polar: se reemplaza la cabeza por una nueva. En la etapa A se congelan el stem y las tres primeras etapas (con su BatchNorm en modo evaluación) y se entrenan la última etapa, el embedding y la cabeza; en la etapa B se descongela toda la red con una tasa de aprendizaje diez veces menor.
  \item \textbf{Alternativa solo Polar}: mismo procedimiento sin preentrenamiento, seleccionado con una bandera de configuración. Permite medir si el preentrenamiento realmente ayuda.
\end{enumerate}

\begin{table}[htbp]
\centering
\caption{Hiperparámetros por etapa de entrenamiento.}
\label{tab:hiperparametros}
\small
\begin{tabularx}{\textwidth}{X c c X}
\toprule
\textbf{Etapa} & \textbf{Épocas} & \textbf{Tasa de aprendizaje} & \textbf{Qué se entrena} \\
\midrule
Preentrenamiento (público) & 30 & $10^{-3}$ & Toda la red y la cabeza \\
Ajuste fino A & 5 & $10^{-3}$ & Última etapa, embedding y cabeza nueva \\
Ajuste fino B & 20 & $10^{-4}$ & Toda la red \\
Solo Polar & 40 & $10^{-3}$ & Toda la red y la cabeza \\
\bottomrule
\end{tabularx}
\end{table}

Se utiliza el optimizador AdamW con $\text{weight decay} = 10^{-4}$, planificador de coseno (\emph{CosineAnnealingLR}), tamaño de lote de 64 y semilla fija igual a 42. La detención temprana se basa en el \textbf{EER de validación} con paciencia de 5 épocas, y no en la exactitud de clasificación, porque el objetivo es separar personas que la red no ha visto. Durante el entrenamiento se aplica aumento de datos: ruido gaussiano (desviación 0.02), deformación de magnitud (0.10), deformación temporal (0.05) y deriva de línea base (amplitud 0.10).

\subsection{Partición de los datos}

Las particiones se hacen siempre \textbf{por persona} y nunca por ventana, porque las ventanas de una misma persona están correlacionadas y mezclarlas inflaría el desempeño.

\begin{table}[htbp]
\centering
\caption{Grupos de personas utilizados.}
\label{tab:particion}
\small
\begin{tabularx}{\textwidth}{l X X}
\toprule
\textbf{Grupo} & \textbf{Uso} & \textbf{Regla} \\
\midrule
Personas de ECG-ID & Preentrenamiento & Nunca se mezclan con las del Polar \\
Polar: entrenamiento & Ajuste fino o entrenamiento solo Polar & Sesiones 1 y 2 entrenan; la sesión 3 sirve como verificación entre días \\
Polar: validación & Detención temprana y umbral & Personas no vistas; nunca actualizan los pesos \\
Polar: \emph{hold-out} & Prueba final & Se usan una sola vez, al final \\
\bottomrule
\end{tabularx}
\end{table}

Con 20 voluntarios la partición prevista es de 14 personas de entrenamiento, 2 de validación y 4 de prueba (\emph{hold-out}), con un mínimo de 2 personas \emph{hold-out}. Como mejora opcional, puede usarse validación cruzada por persona en 5 pliegues de 4 personas, de modo que las 20 sirvan una vez como personas no vistas. La lista de validación y de \emph{hold-out} se elige al azar una sola vez, después de recolectar los datos y antes de entrenar con datos Polar, y se registra en el archivo de configuración; el código se detiene con un error si alguna de esas personas aparece en los datos de entrenamiento. Con 4 personas de prueba hay solo 12 combinaciones plantilla--impostor, por lo que el EER final tendrá intervalos de confianza amplios.

\subsection{Métricas de evaluación}

\begin{itemize}
  \item \textbf{Pares genuinos}: ventanas de una sesión posterior contra la plantilla de la misma persona (entre sesiones).
  \item \textbf{Pares impostores}: ventanas de una persona contra la plantilla de otra.
  \item \textbf{FAR} (tasa de falsa aceptación) y \textbf{FRR} (tasa de falso rechazo) para un umbral dado, y la \textbf{EER}, el punto en que ambas son iguales:
  \begin{equation}
    \mathrm{FAR}(\tau) = \mathrm{FRR}(\tau) = \mathrm{EER}
  \end{equation}
  \item Curva ROC y AUC.
  \item El umbral $\tau$ se elige en el punto de EER con datos de validación únicamente. Los intervalos de confianza se obtienen por \emph{bootstrap} remuestreando personas (1000 repeticiones), y cada cifra se reporta junto con el número de pares genuinos e impostores en que se basa.
\end{itemize}

\subsection{Diseño de la decisión continua}
\label{sec:decision}

La decisión final no se basa en una sola ventana, pues una lectura aislada puede fallar por un movimiento o por mal contacto de los electrodos. La regla de decisión continua es la siguiente:

\begin{enumerate}
  \item Cada 2.5~s se obtiene una ventana nueva y se calcula su similitud coseno con la plantilla del usuario.
  \item Si la similitud es mayor o igual que el umbral $\tau$, la ventana se acepta y el contador de rechazos consecutivos se reinicia a cero.
  \item Si es menor, se rechaza y el contador aumenta en uno. Mientras el contador sea menor que 3, el usuario no percibe nada.
  \item Con 3 rechazos consecutivos aparece una ventana emergente que pide al usuario sentarse derecho y permanecer quieto unos segundos, y el sistema toma una nueva muestra.
  \item Si la nueva muestra es aceptada, el examen continúa y el contador se reinicia. Si es rechazada, el examen se bloquea.
\end{enumerate}

Dado que la primera ventana completa se obtiene a los 5~s y las siguientes a intervalos de 2.5~s, el aviso al usuario aparece como mínimo unos 10~s después de que comienza una suplantación, y el bloqueo ocurre después de la nueva muestra, de otros 5~s como mínimo. Si cada rechazo injustificado ocurriera de forma independiente con probabilidad igual a la FRR, la probabilidad de 3 consecutivos sería $\mathrm{FRR}^3$; esto es solo una aproximación, pues las ventanas traslapadas no son independientes, y la tasa real de bloqueos injustificados se medirá con grabaciones completas en el Capítulo~\ref{cap:resultados}. El valor de 3 rechazos es un parámetro de diseño que puede ajustarse con los resultados. Esta lógica es independiente de la red, que solo produce los embeddings.

\subsection{Simulación y pruebas con señales sintéticas}

Antes de contar con datos reales se construyó un generador de señales ECG sintéticas para verificar el funcionamiento de cada módulo. Con ellas se comprobó que la red devuelve la forma de salida esperada $(B, 128)$ con norma igual a 1, que el resultado es determinista en modo evaluación, y que un entrenamiento pequeño sobre un lote reduce la pérdida. También se comprobó que la evaluación da un EER cercano a 0 con puntuaciones perfectamente separables y cercano a 0.5 con puntuaciones aleatorias, que un guardado y carga del modelo producen los mismos embeddings, y que el programa se detiene si una persona \emph{hold-out} aparece en el entrenamiento.

\section{Construcción e implementación del prototipo}
\label{sec:construccion}

Esta sección describe cómo el diseño se convirtió en un prototipo funcional: la organización del software, la adquisición de datos con el sensor, el protocolo de recolección, las pruebas parciales y los cambios realizados durante las iteraciones.

\subsection{Hardware}

El prototipo utiliza una banda pectoral Polar H10, que registra un canal de ECG a 130~Hz y lo transmite por Bluetooth Low Energy a una computadora. Antes de cada grabación se humedecen los electrodos de la banda para asegurar el contacto. El entrenamiento se realiza en una computadora portátil HP OmniBook X con 16~GB de RAM y CPU de 6 núcleos, sin GPU. \pendiente{agregar fotografía de la banda}

\subsection{Implementación del software}

El software se desarrolló en Python 3.14 con PyTorch 2.14 (versión para CPU), NumPy 2.3, SciPy 1.16, scikit-learn, Matplotlib 3.10 y Bleak 3.0 para la conexión Bluetooth, y se organizó de forma modular para facilitar su explicación, prueba y mantenimiento.

\begin{table}[htbp]
\centering
\caption{Módulos del software.}
\label{tab:modulos}
\small
\begin{tabularx}{\textwidth}{l X}
\toprule
\textbf{Módulo} & \textbf{Función} \\
\midrule
\texttt{config} & Todos los parámetros, la semilla y la lista de personas \emph{hold-out} \\
\texttt{record}, \texttt{check} & Conexión BLE con el Polar y chequeo rápido de la lectura antes de grabar \\
\texttt{protocol}, \texttt{session}, \texttt{status} & Protocolo de sesiones, grabación guiada con nombres automáticos de archivo y seguimiento del avance \\
\texttt{quality} & Control de calidad de la señal (señal plana, ruido, picos anómalos, número de latidos) \\
\texttt{windowing}, \texttt{data} & Filtrado, ventaneo, conjuntos de datos, particiones por persona y aumento de datos \\
\texttt{public\_data} & Carga y remuestreo de la base pública ECG-ID \\
\texttt{model} & Red neuronal de embeddings \\
\texttt{train}, \texttt{checkpoint} & Preentrenamiento, ajuste fino, entrenamiento solo Polar y guardado de modelos con archivo JSON de configuración y métricas \\
\texttt{evaluate} & EER, ROC, FAR/FRR, intervalos \emph{bootstrap} y tablas de resultados \\
\texttt{synthetic} & Señales sintéticas para pruebas \\
\bottomrule
\end{tabularx}
\end{table}

Cada etapa de entrenamiento se ejecuta con un solo comando y los modelos se guardan como versiones numeradas junto con un archivo que registra la configuración, la semilla, la lista de personas \emph{hold-out}, la versión de los datos y las métricas. Solo se almacenan los pesos del codificador.

\subsection{Protocolo de recolección de datos}

Con el Polar H10 se grabaron 20 voluntarios con identificadores anónimos fijos (P01, P02, \dots). Cada voluntario realiza \textbf{3 sesiones en días distintos}, con una separación de 2 a 3 días, para medir la variación de la señal entre días. Cada sesión dura aproximadamente 5.5 minutos y se divide en cuatro actividades:

\begin{table}[htbp]
\centering
\caption{Actividades de cada sesión de grabación.}
\label{tab:sesion}
\small
\begin{tabular}{l c}
\toprule
\textbf{Actividad} & \textbf{Duración (min)} \\
\midrule
Reposo & 2 \\
Lectura & 1.5 \\
Resolución de problemas & 1 \\
Hablando & 0.5 \\
\bottomrule
\end{tabular}
\end{table}

El programa guía al voluntario, evalúa la calidad de cada actividad y permite repetir una si salió mal. Los datos crudos no se suben al repositorio. Cada participante firmó un formato de consentimiento informado que especificaba las actividades que se realizarían, los cuidados durante la grabación, el uso que se daría a los datos (exclusivamente de investigación, con identificadores anónimos) y que la grabación no implicaba riesgos para su salud. \pendiente{elaborar y recabar el formato firmado; anexarlo al documento}

\subsection{Calibración y validación de instrumentos}

La calidad de cada actividad se verifica automáticamente con \texttt{quality}, y antes de grabar se corre un chequeo de unos 12~s con \texttt{check} que advierte si hay que ajustar la banda. Para el preentrenamiento se comprobó la frecuencia de muestreo de la base pública y se remuestreó a 130~Hz.

\subsection{Integración y pruebas parciales}

Cada módulo cuenta con pruebas automáticas con señales sintéticas (por ejemplo, forma de salida y norma del embedding, ausencia de fugas de personas \emph{hold-out} hacia el entrenamiento, evaluación con casos de EER conocido, guardado y carga de modelos, y comprobación de la calidad de señal). Primero se verificó el modelo y los datos por separado, después el entrenamiento y la evaluación, y finalmente la adquisición con el sensor.

\subsection{Iteraciones del prototipo}

\begin{table}[htbp]
\centering
\caption{Iteraciones del prototipo.}
\label{tab:iteraciones}
\small
\begin{tabularx}{\textwidth}{c X X}
\toprule
\textbf{Versión} & \textbf{Implementación principal} & \textbf{Propósito} \\
\midrule
1 & Red de embeddings, entrenamiento y evaluación; preentrenamiento con ECG-ID (modelo v1) & Comprobar que la red aprende a distinguir personas \\
2 & Adquisición con el Polar, control de calidad y sesiones guiadas de cuatro actividades & Obtener datos propios con un protocolo repetible \\
3 & Recolección con los 20 voluntarios y ajuste fino con el Polar & Adaptar la red al sensor y a las personas del estudio \\
4 & Lógica de decisión continua (3 rechazos, nueva muestra, bloqueo) y, de ser necesario, mejoras de preprocesamiento & Detectar un impostor durante el examen \\
5 & Interfaz del examen de unas 30 preguntas de comprensión lectora y pruebas de tiempo de detección & Integrar el sistema completo \\
\bottomrule
\end{tabularx}
\end{table}

\pendiente{las versiones 1 y 2 corresponden a lo ya implementado (confirmado en el historial de git); las versiones 3 a 5 son el plan de trabajo actual. Ajustar al entregar, indicando para cada versión el problema detectado, el cambio realizado y la razón técnica}

\subsection{Simulación del examen}

El examen simulado consta de unas 30 preguntas de comprensión lectora: se presentan lecturas y, después de cada una, aproximadamente 5 preguntas. Mientras el usuario lo contesta, el sistema verifica su identidad de manera continua con la lógica descrita en la Sección~\ref{sec:decision}. \pendiente{definir la plataforma, la duración y las lecturas empleadas}

\section{Presentación de la solución desarrollada}
\label{sec:solucion}

Esta sección presenta el resultado final del proceso de diseño y construcción: qué solución se obtuvo, cómo funciona y cuáles son sus características. Aquí se describe lo que se construyó; el Capítulo~\ref{cap:resultados} demuestra qué tan bien funcionó.

\subsection{Descripción de la solución}

La solución desarrollada es un sistema de autenticación continua que verifica la identidad de una persona durante un examen simulado a partir de su ECG. El usuario se registra una vez mediante una sesión de enrolamiento, de la cual se obtiene su plantilla. Durante el examen, la banda Polar H10 transmite la señal por Bluetooth; el sistema la filtra y la divide en ventanas de 5~s con 50\,\% de traslape, y una red neuronal convolucional 1D residual genera un embedding por ventana. La similitud coseno entre el embedding y la plantilla se compara con un umbral calibrado con datos de validación para aceptar o rechazar la identidad. Los rechazos aislados no se muestran al usuario; con 3 rechazos consecutivos se le pide una nueva muestra y, si esta también falla, el examen se bloquea (Sección~\ref{sec:decision}). La primera decisión está disponible a los 5~s de iniciada la adquisición y las siguientes cada 2.5~s. La interfaz del examen está por definirse.

\subsection{Secuencia de funcionamiento}

\begin{enumerate}
  \item Se coloca la banda y se verifica la calidad de la señal.
  \item El usuario realiza el registro y se calcula su plantilla.
  \item Inicia el examen y comienza la adquisición continua.
  \item Cada 2.5~s (por el traslape) se procesa una nueva ventana de 5~s.
  \item Se calcula la similitud con la plantilla y se decide si se acepta la identidad.
  \item Los rechazos aislados no se notifican. Con 3 consecutivos aparece un aviso para sentarse derecho y se toma una nueva muestra; si también es rechazada, el examen se bloquea.
\end{enumerate}

\subsection{Ficha técnica}

\begin{table}[htbp]
\centering
\caption{Ficha técnica de la solución.}
\label{tab:ficha}
\small
\begin{tabularx}{\textwidth}{l X}
\toprule
\textbf{Característica} & \textbf{Descripción} \\
\midrule
Tipo de sistema & Autenticación biométrica continua basada en ECG \\
Sensor & Polar H10, un canal, 130~Hz, conexión BLE \\
Entrada del modelo & Ventana de 650 muestras (5~s), filtrada y normalizada \\
Modelo & CNN 1D residual de 371{,}024 parámetros \\
Salida & Embedding de 128 dimensiones con norma L2 igual a 1 \\
Comparación & Similitud coseno con la plantilla del usuario \\
Entrenamiento & Preentrenamiento con ECG-ID y ajuste fino con voluntarios del Polar \\
Hardware de cómputo & CPU, sin GPU \\
Lenguaje y bibliotecas & Python, PyTorch, NumPy, SciPy, scikit-learn \\
Datos registrados & Señal por sesión con ID anónimo, calidad y metadatos \\
Aplicación & Verificación continua de identidad en un examen simulado \\
\bottomrule
\end{tabularx}
\end{table}

\subsection{Relación entre requerimientos e implementación}

\begin{table}[htbp]
\centering
\caption{Relación general entre requerimientos e implementación.}
\label{tab:req_impl}
\small
\begin{tabularx}{\textwidth}{l X}
\toprule
\textbf{Requerimiento} & \textbf{Implementación en la solución} \\
\midrule
RF-01 Adquirir ECG & Módulo \texttt{record} con el Polar H10 por BLE \\
RF-02 Calidad de la señal & Módulos \texttt{quality} y \texttt{check} \\
RF-03 Embedding por ventana & Red 1D-CNN residual del módulo \texttt{model} \\
RF-04 Plantilla del usuario & Promedio de embeddings de la sesión de enrolamiento \\
RF-05 Decisión por similitud & Similitud coseno y umbral elegido en validación \\
RD-01, RD-02 Desempeño en personas nuevas y entre días & Partición por persona y evaluación con \texttt{evaluate} \\
RF-06, RD-03, RD-04 Detección de impostor sin molestar al legítimo & Regla de 3 rechazos consecutivos, nueva muestra y bloqueo \\
RNF-01 Sin GPU & Red de 371{,}024 parámetros entrenable en CPU \\
RNF-02 Reproducibilidad & Semilla fija y configuración centralizada en \texttt{config} \\
RNF-03 Módulos separados & Red, plantillas y decisión en componentes independientes \\
RNF-04 Anonimato & IDs fijos y datos fuera del repositorio \\
\bottomrule
\end{tabularx}
\end{table}

\subsection{Limitaciones conocidas}

\begin{itemize}
  \item El número de voluntarios es reducido, por lo que el EER del conjunto \emph{hold-out} se apoya en pocas combinaciones de personas e intervalos de confianza amplios. Los resultados se consideran una prueba de concepto.
  \item La base pública ECG-ID usa derivaciones y equipo distintos al de la banda Polar, lo que genera una diferencia de dominio que el ajuste fino busca reducir.
  \item Se utiliza un solo canal de ECG y un solo modelo de sensor.
  \item Las sesiones se grabaron en un periodo corto de semanas, por lo que no se evalúan cambios de largo plazo en la señal.
  \item La actividad durante el examen simulado es principalmente sedentaria; no se evaluó movimiento intenso ni estados emocionales extremos.
  \item Las funciones de la pérdida de tripleta y ArcFace no se implementaron en esta versión.
\end{itemize}


\section*{Cierre del capítulo}

En este capítulo se definieron los requerimientos funcionales, de desempeño y no funcionales del sistema; se generaron y compararon alternativas mediante técnicas de creatividad y una matriz de decisión; y se seleccionó una arquitectura basada en una red neuronal convolucional unidimensional con bloques residuales que produce embeddings de identidad. También se describieron el diseño de la red, el preprocesamiento, la estrategia de entrenamiento y la partición de datos por persona, así como la implementación modular del prototipo y su protocolo de recolección. En el siguiente capítulo se presentan los resultados de las pruebas experimentales, el análisis de los indicadores de desempeño y la comprobación del cumplimiento de los objetivos y de la hipótesis.
 desde el documento principal.
% ============================================================

% Marca visible para lo que el equipo debe completar o confirmar.
\newcommand{\pendiente}[1]{\textbf{\textcolor{red}{[POR DEFINIR: #1]}}}

\chapter{Diseño y desarrollo de la solución}
\label{cap:diseno}

En este capítulo se describe el proceso de diseño y desarrollo de un sistema de autenticación continua basado en la señal electrocardiográfica (ECG) adquirida con una banda Polar H10. Se presentan los requerimientos de la solución, las alternativas tecnológicas consideradas, el diseño detallado del sistema, la construcción del prototipo y las características de la solución finalmente desarrollada.

El proyecto se llevó a cabo mediante una metodología de prototipado iterativo. Primero se construyó un extractor de embeddings entrenado con una base de datos pública, después se incorporó la adquisición de señal con el sensor y el protocolo de recolección con voluntarios, y finalmente se integraron los módulos de decisión continua. Cada iteración permitió identificar problemas y mejorar el sistema. Las tablas de resultados, el desempeño medido y la comprobación de la hipótesis se presentan en el Capítulo~\ref{cap:resultados}.

\section{Especificación de requerimientos}
\label{sec:requerimientos}

Esta sección establece qué debe hacer el sistema y bajo qué condiciones debe funcionar. Los requerimientos sirven como referencia para diseñar la solución y para comprobar, en el Capítulo~\ref{cap:resultados}, si el prototipo cumple con lo solicitado.

\subsection{Descripción general del sistema}

El sistema simula un examen en computadora. Mientras la persona lo contesta, una banda Polar H10 colocada en el pecho registra su ECG. Una red neuronal convierte ventanas sucesivas de la señal en vectores de características (\emph{embeddings}) y los compara con la plantilla creada en el registro inicial del usuario. De esta forma el sistema verifica de manera continua que quien contesta el examen es la misma persona que se registró, y no solo en el momento de iniciar sesión.

\subsection{Requerimientos funcionales}

\begin{itemize}
  \item Adquirir la señal ECG del Polar H10 por Bluetooth Low Energy a 130~Hz.
  \item Verificar la calidad de la señal antes y durante la grabación (señal plana, ruido excesivo, picos anómalos, latidos insuficientes).
  \item Preprocesar la señal: filtro pasa-banda de 0.5 a 40~Hz, segmentación en ventanas de 5~s con 50\,\% de traslape y normalización z-score por ventana.
  \item Generar un embedding L2-normalizado de 128 dimensiones por ventana.
  \item Crear la plantilla de un usuario durante el registro (sesión de enrolamiento) a partir del promedio de sus embeddings.
  \item Calcular la similitud coseno entre cada ventana nueva y la plantilla, y decidir si la identidad se acepta o se rechaza.
  \item Aplicar una regla de decisión continua: los rechazos aislados no se muestran al usuario; tras 3 rechazos consecutivos se le pide una nueva muestra (por ejemplo, sentarse derecho y permanecer quieto unos segundos) y, si esta tampoco es aceptada, se bloquea el examen.
  \item Entrenar y evaluar el modelo con personas que no participaron en el entrenamiento.
  \item Registrar los datos de cada sesión con identificadores anónimos, para poder repetir los análisis.
\end{itemize}

\subsection{Requerimientos no funcionales y restricciones}

\begin{itemize}
  \item El entrenamiento debe ejecutarse en una computadora sin GPU (16~GB de RAM, CPU de 6 núcleos), en un tiempo del orden de minutos a una hora.
  \item Los experimentos deben ser reproducibles: semilla fija, un único archivo de configuración y un comando por etapa de entrenamiento.
  \item La red neuronal debe permanecer separada de la base de plantillas y de la lógica de decisión, de modo que cada componente pueda sustituirse sin modificar los demás.
  \item Los datos fisiológicos deben tratarse de forma anónima y no deben publicarse en el repositorio.
  \item El modelo debe poder explicarse por completo en cada una de sus partes por los integrantes del equipo.
  \item Restricciones del proyecto: una muestra de 20 voluntarios, un solo canal de ECG, un solo sensor disponible y un plazo de desarrollo que concluye en noviembre de 2026.
\end{itemize}

\subsection{Tabla de requerimientos}

\begin{table}[htbp]
\centering
\caption{Especificación de requerimientos y criterios de aceptación.}
\label{tab:requerimientos}
\small
\begin{tabularx}{\textwidth}{l X l X}
\toprule
\textbf{Código} & \textbf{Requerimiento} & \textbf{Tipo} & \textbf{Criterio de aceptación} \\
\midrule
RF-01 & Adquirir ECG del Polar H10 a 130~Hz & Funcional & Se guardan sesiones completas sin pérdida de conexión \\
RF-02 & Evaluar la calidad de la señal & Funcional & Una sesión con señal defectuosa se detecta y se repite \\
RF-03 & Generar un embedding por ventana de 5~s & Funcional & Salida de dimensión 128 con norma L2 igual a 1 \\
RF-04 & Crear la plantilla de cada usuario & Funcional & Plantilla construida solo con la sesión de enrolamiento \\
RF-05 & Decidir identidad por similitud coseno & Funcional & Umbral fijado solo con datos de validación \\
RF-06 & Aplicar la regla de decisión continua & Funcional & Rechazos aislados ocultos; con 3 consecutivos se pide nueva muestra y, si falla, se bloquea \\
RD-01 & Distinguir personas nunca vistas por la red & Desempeño & EER medido en personas \emph{hold-out}, claramente por debajo del azar (50\,\%) y con intervalo de confianza que no lo incluya \\
RD-02 & Mantener el desempeño entre días & Desempeño & EER entre sesiones de días distintos reportado con intervalo de confianza \\
RD-03 & Detectar a un impostor durante el examen & Desempeño & Tras 3 rechazos consecutivos (unos 10~s) se solicita una nueva muestra; si esta también falla, se bloquea el examen \\
RD-04 & No interrumpir innecesariamente al usuario legítimo & Desempeño & Los rechazos aislados no se notifican; la tasa de bloqueos injustificados se reporta en el Capítulo~\ref{cap:resultados} \\
RNF-01 & Entrenar sin GPU & No funcional & Entrenamiento completo en CPU en menos de una hora \\
RNF-02 & Reproducibilidad & No funcional & Misma semilla y configuración producen el mismo modelo \\
RNF-03 & Separación entre red, plantillas y decisión & No funcional & Cada módulo se prueba de forma independiente \\
RNF-04 & Anonimato de los participantes & No funcional & Identificadores fijos tipo P01 y datos excluidos del repositorio \\
\bottomrule
\end{tabularx}
\end{table}

No se fija una meta numérica rígida de EER: con 20 participantes y un solo canal, el objetivo es demostrar la viabilidad del enfoque. El criterio es que el EER quede claramente por debajo del azar y se compare con los valores reportados en el estado del arte del Capítulo~\ref{cap:marco}. Por ello, los resultados del proyecto se consideran una prueba de concepto.

\section{Generación, evaluación y selección de alternativas}
\label{sec:alternativas}

Esta sección muestra cómo se generaron distintas propuestas para resolver el problema y por qué se eligió una de ellas. Las técnicas de creatividad no son las alternativas finales: son procedimientos para generar ideas que después se organizan, comparan y seleccionan.

% NOTA PARA EL EQUIPO: las ideas, tablas y calificaciones de esta sección son una
% PROPUESTA de partida. Deben discutirse y ajustarse según lo que el equipo realmente
% consideró, y sustentarse con literatura o pruebas preliminares. Conserven evidencias
% (fotos de la lluvia de ideas, mapa mental) para anexarlas.

\subsection{Técnicas para la generación de ideas}

\subsubsection{Lluvia de ideas}

Pregunta generadora: ¿qué sensor, qué representación de la señal y qué modelo podrían utilizarse para verificar de forma continua la identidad de una persona a partir de su ECG? Las ideas obtenidas fueron: banda Polar H10, un sensor de ECG de diseño propio, otros equipos comerciales de mayor costo (por ejemplo, Movesense o Zephyr BioHarness), señal cruda o latidos segmentados, rasgos manuales (intervalos y amplitudes de las ondas P, QRS y T), máquina de vectores de soporte (SVM), red neuronal convolucional 1D como clasificador, red 1D con embeddings, pérdida de tripleta, preentrenamiento con una base de datos pública y suavizado de las decisiones en el tiempo. \pendiente{confirmar con el equipo qué equipo comercial más caro se consideró}

\subsubsection{Mapa mental}

Se colocó en el centro el concepto <<autenticación continua con ECG>> y se crearon ramas para adquisición, preprocesamiento, representación de la señal, modelo, entrenamiento, decisión continua, interfaz y validación. El mapa permitió verificar que las propuestas contemplaran el sistema completo y no solamente la red neuronal. \pendiente{insertar figura del mapa mental}

\subsubsection{SCAMPER}

\begin{table}[htbp]
\centering
\caption{Aplicación de SCAMPER al sistema de autenticación continua.}
\label{tab:scamper}
\small
\begin{tabularx}{\textwidth}{l X X}
\toprule
\textbf{Acción} & \textbf{Pregunta aplicada} & \textbf{Ejemplo de idea} \\
\midrule
Sustituir & ¿Qué componente puede reemplazarse? & Sustituir rasgos manuales por rasgos aprendidos por una red \\
Combinar & ¿Qué elementos pueden trabajar juntos? & Combinar preentrenamiento público con ajuste fino con datos del Polar \\
Adaptar & ¿Qué solución externa puede aprovecharse? & Adaptar el aprendizaje de métricas usado en reconocimiento facial \\
Modificar & ¿Qué característica puede cambiarse? & Cambiar la longitud de la ventana de 5~s \\
Otros usos & ¿Puede un elemento cumplir otra función? & Usar la penúltima capa de un clasificador como huella de identidad \\
Eliminar & ¿Qué elemento no es indispensable? & Eliminar la detección explícita de picos R y trabajar sobre la señal filtrada \\
Reordenar & ¿Qué ocurre al cambiar el orden? & Suavizar decisiones en el tiempo en lugar de decidir por ventana aislada \\
\bottomrule
\end{tabularx}
\end{table}

\subsubsection{Combinación de elementos y matriz morfológica}

La matriz morfológica combina opciones de distintos subsistemas para construir alternativas completas y técnicamente coherentes.

\begin{table}[htbp]
\centering
\caption{Matriz morfológica del sistema.}
\label{tab:morfologica}
\small
\begin{tabularx}{\textwidth}{l X X X}
\toprule
\textbf{Función} & \textbf{Opción 1} & \textbf{Opción 2} & \textbf{Opción 3} \\
\midrule
Adquisición & Polar H10 & Sensor de ECG de diseño propio & Equipo comercial de mayor costo \\
Entrada del modelo & Rasgos manuales & Ventana cruda filtrada & Latidos segmentados \\
Modelo & SVM & CNN 1D como clasificador & CNN 1D con embeddings \\
Pérdida & Entropía cruzada & Tripleta & ArcFace \\
Datos de entrenamiento & Solo Polar & Preentrenado y ajustado & Solo base pública \\
Decisión & Por ventana & Promedio móvil & Acumulación de evidencia \\
\bottomrule
\end{tabularx}
\end{table}

\subsection{Alternativas de solución obtenidas}

\begin{table}[htbp]
\centering
\caption{Alternativas conceptuales.}
\label{tab:alternativas}
\small
\begin{tabularx}{\textwidth}{l X X X}
\toprule
\textbf{Alternativa} & \textbf{Configuración} & \textbf{Ventaja principal} & \textbf{Limitación principal} \\
\midrule
A: clásica & Polar H10, rasgos manuales, SVM, decisión por ventana & Pocos datos necesarios y fácil de explicar & Depende de detectar con precisión las ondas; menor capacidad de generalizar a personas nuevas \\
B: embeddings con preentrenamiento & Polar H10, ventana filtrada, CNN 1D residual con embeddings, preentrenamiento público y ajuste fino, entropía cruzada & Aprende rasgos automáticamente; permite registrar personas nuevas sin reentrenar; el preentrenamiento compensa pocos voluntarios & Diferencia entre la base pública y el Polar; requiere más tiempo de entrenamiento \\
C: aprendizaje de métricas avanzado & Igual que B pero con pérdida de tripleta o ArcFace y detección de picos R & Mayor separación potencial entre identidades & Mayor complejidad y sensibilidad a hiperparámetros con pocos datos \\
\bottomrule
\end{tabularx}
\end{table}

\subsection{Criterios de evaluación}

\begin{table}[htbp]
\centering
\caption{Criterios de evaluación y ponderación.}
\label{tab:criterios}
\small
\begin{tabularx}{\textwidth}{l l X}
\toprule
\textbf{Criterio} & \textbf{Ponderación} & \textbf{Justificación} \\
\midrule
Capacidad de generalizar a personas nuevas & 30\,\% & El sistema debe funcionar con usuarios que la red nunca vio \\
Viabilidad con pocos datos & 20\,\% & Solo se cuenta con 20 voluntarios, una cantidad pequeña para entrenar redes profundas \\
Costo computacional & 15\,\% & Solo se dispone de CPU \\
Facilidad de implementación y explicación & 20\,\% & El equipo debe dominar cada componente y cumplir el plazo \\
Aprovechamiento del sensor disponible & 15\,\% & El Polar H10 ya está en poder del equipo \\
\bottomrule
\end{tabularx}
\end{table}

Cada alternativa se califica del 1 al 5, donde 1 es un desempeño muy deficiente y 5 uno excelente. El resultado ponderado se obtiene multiplicando la calificación por la ponderación; la suma permite comparar las alternativas.

\begin{table}[htbp]
\centering
\caption{Matriz de decisión (calificaciones preliminares, por justificar).}
\label{tab:decision}
\small
\begin{tabular}{l c c c c}
\toprule
\textbf{Criterio} & \textbf{Pond.} & \textbf{A} & \textbf{B} & \textbf{C} \\
\midrule
Generalización a personas nuevas & 0.30 & 2 & 4 & 5 \\
Viabilidad con pocos datos & 0.20 & 4 & 4 & 2 \\
Costo computacional & 0.15 & 5 & 4 & 3 \\
Facilidad de implementación y explicación & 0.20 & 4 & 4 & 2 \\
Aprovechamiento del sensor & 0.15 & 4 & 4 & 4 \\
\midrule
\textbf{Total ponderado} & 1.00 & 3.55 & 4.00 & 3.35 \\
\bottomrule
\end{tabular}
\end{table}

\pendiente{justificar cada calificación con literatura, pruebas preliminares o cálculos; las cifras mostradas son un borrador}

\subsection{Selección y justificación}

Después de aplicar las técnicas de creatividad se obtuvieron tres alternativas de solución, evaluadas según su capacidad de generalizar a personas nuevas, su viabilidad con pocos datos, su costo computacional, su facilidad de implementación y el aprovechamiento del sensor disponible. La alternativa B obtuvo la mayor puntuación porque ofrece un equilibrio adecuado entre capacidad de generalización y viabilidad: el uso de embeddings permite registrar usuarios nuevos solo calculando su plantilla, y el preentrenamiento con una base de datos pública mitiga la escasez de voluntarios. La alternativa C se consideró una mejora futura en caso de que el error obtenido con B no fuera suficiente, por lo que la pérdida de tripleta y ArcFace se dejaron deliberadamente fuera de esta versión. En cuanto al sensor, se consideró construir uno propio, pero se encontró que el Polar H10 es más preciso, por lo que se descartó. Los equipos comerciales de mayor costo no ofrecían una ventaja que justificara su precio, y el Polar H10 cuenta con estudios que lo validan frente a un ECG de referencia. Por lo anterior, se seleccionó la alternativa B, con el Polar H10, como base para el diseño detallado.

\section{Modelado, cálculos, simulación y diseño detallado}
\label{sec:diseno}

Esta sección explica técnicamente cómo se diseñó la alternativa seleccionada antes de construirla, de modo que otra persona pueda comprender la arquitectura del sistema y reproducir el diseño.

\subsection{Arquitectura general}

El sistema se organiza como una cadena de bloques independientes:
\begin{center}
\emph{Polar H10} $\rightarrow$ adquisición BLE $\rightarrow$ control de calidad $\rightarrow$ filtrado y ventaneo $\rightarrow$ red de embeddings $\rightarrow$ similitud coseno con la plantilla $\rightarrow$ decisión
\end{center}
\pendiente{insertar diagrama de bloques}

La red (los pesos) se mantiene separada de la base de plantillas y de la lógica de decisión. Así, registrar a un usuario nuevo solo requiere calcular su plantilla, sin reentrenar la red.

\subsection{Adquisición y preprocesamiento de la señal}

La banda Polar H10 entrega el ECG a $f_s = 130$~Hz. La señal se procesa de la siguiente forma:
\begin{enumerate}
  \item \textbf{Filtro pasa-banda} de 0.5 a 40~Hz, que elimina la deriva de la línea base y atenúa el ruido muscular de alta frecuencia. La frecuencia de Nyquist es de 65~Hz, por lo que la banda es válida.
  \item \textbf{Ventaneo} en segmentos de 5~s, es decir, $5 \times 130 = 650$ muestras, con 50\,\% de traslape.
  \item \textbf{Normalización z-score por ventana}, para eliminar diferencias de amplitud debidas al contacto de los electrodos:
  \begin{equation}
    \tilde{x}_i = \frac{x_i - \mu}{\sigma}
  \end{equation}
  donde $\mu$ y $\sigma$ son la media y la desviación estándar de la ventana.
\end{enumerate}
Cada ventana es un tensor de forma $(1, 650)$. El sistema no detecta explícitamente los picos R para el modelo: la red extrae por sí misma las características morfológicas y temporales.

\subsection{Red neuronal de embeddings}

\subsubsection{Tipo de red}

Se utiliza una red neuronal convolucional unidimensional (1D-CNN) con bloques residuales, usada como \textbf{extractor de embeddings} (aprendizaje de métricas) y no como clasificador final. Su salida es un vector de 128 dimensiones con norma L2 igual a 1, de modo que la similitud coseno entre dos ventanas indica si provienen de la misma persona.

\subsubsection{Arquitectura capa por capa}

\begin{table}[htbp]
\centering
\caption{Arquitectura de la red para una ventana de entrada $(1, 650)$.}
\label{tab:arquitectura}
\small
\begin{tabularx}{\textwidth}{c l X c c c}
\toprule
\textbf{\#} & \textbf{Bloque} & \textbf{Operación} & \textbf{Kernel / stride} & \textbf{Canales} & \textbf{Salida} \\
\midrule
0 & Entrada & Ventana z-score & -- & 1 & $1 \times 650$ \\
1 & Stem & Conv1d + BatchNorm + ReLU & 7 / 1 & $1 \rightarrow 16$ & $16 \times 650$ \\
2 & Etapa 1 & Bloque residual & 5 / 2 & $16 \rightarrow 32$ & $32 \times 325$ \\
3 & Etapa 2 & Bloque residual & 5 / 2 & $32 \rightarrow 64$ & $64 \times 163$ \\
4 & Etapa 3 & Bloque residual & 5 / 2 & $64 \rightarrow 128$ & $128 \times 82$ \\
5 & Etapa 4 & Bloque residual & 5 / 2 & $128 \rightarrow 128$ & $128 \times 41$ \\
6 & Pooling global & Promedio sobre el tiempo & -- & 128 & 128 \\
7 & Embedding & Capa densa & -- & $128 \rightarrow 128$ & 128 \\
8 & Normalización & L2 & -- & 128 & 128 (norma 1) \\
9 & Cabeza & Capa densa sobre el embedding sin normalizar (solo entrenamiento) & -- & $128 \rightarrow N$ & $N$ \\
\bottomrule
\end{tabularx}
\end{table}

Todas las convoluciones usan \texttt{bias=False}, pues la normalización por lotes (BatchNorm) ya aporta el desplazamiento. La red tiene \textbf{371{,}024 parámetros} en el codificador; la cabeza agrega $128N + N$ parámetros, con $N$ igual al número de personas de entrenamiento, y se descarta al terminar.

\subsubsection{Bloque residual}

Cada etapa usa el mismo bloque. La rama principal aplica una convolución con $k=5$ y \emph{stride} 2, BatchNorm y ReLU, y después una segunda convolución con $k=5$ y \emph{stride} 1 seguida de BatchNorm. El atajo (\emph{skip connection}) es una convolución de $1\times1$ con \emph{stride} 2 seguida de BatchNorm, para que las dimensiones coincidan en la suma. La ReLU se aplica después de sumar ambas ramas:
\begin{equation}
  \mathbf{y} = \mathrm{ReLU}\big(F(\mathbf{x}) + S(\mathbf{x})\big)
\end{equation}
donde $F$ es la rama principal y $S$ el atajo. El \emph{stride} 2 reduce a la mitad la longitud temporal en cada etapa ($650 \rightarrow 325 \rightarrow 163 \rightarrow 82 \rightarrow 41$); con ello el campo receptivo de la última etapa abarca varios latidos.

\subsubsection{Embedding y similitud}

El pooling global promedio hace que la huella sea independiente de en qué punto de la ventana cae cada latido. La salida se normaliza:
\begin{equation}
  \mathbf{e} = \frac{\mathbf{z}}{\lVert \mathbf{z} \rVert_2}
\end{equation}
y la similitud entre una ventana con embedding $\mathbf{e}$ y una plantilla $\mathbf{t}$ es el coseno del ángulo entre ambos, que al estar normalizados equivale al producto punto:
\begin{equation}
  s(\mathbf{e}, \mathbf{t}) = \mathbf{e} \cdot \mathbf{t}
\end{equation}
La plantilla de una persona es el promedio de los embeddings de su sesión de enrolamiento. La identidad se acepta si $s$ supera un umbral $\tau$.

\subsection{Estrategia de entrenamiento}

\begin{enumerate}
  \item \textbf{Preentrenamiento} con la base pública ECG-ID, remuestreada a 130~Hz: la red se entrena como clasificador de identidades con entropía cruzada y \emph{label smoothing} de 0.1. Su propósito es aprender rasgos generales del ECG.
  \item \textbf{Ajuste fino} con los voluntarios del Polar: se reemplaza la cabeza por una nueva. En la etapa A se congelan el stem y las tres primeras etapas (con su BatchNorm en modo evaluación) y se entrenan la última etapa, el embedding y la cabeza; en la etapa B se descongela toda la red con una tasa de aprendizaje diez veces menor.
  \item \textbf{Alternativa solo Polar}: mismo procedimiento sin preentrenamiento, seleccionado con una bandera de configuración. Permite medir si el preentrenamiento realmente ayuda.
\end{enumerate}

\begin{table}[htbp]
\centering
\caption{Hiperparámetros por etapa de entrenamiento.}
\label{tab:hiperparametros}
\small
\begin{tabularx}{\textwidth}{X c c X}
\toprule
\textbf{Etapa} & \textbf{Épocas} & \textbf{Tasa de aprendizaje} & \textbf{Qué se entrena} \\
\midrule
Preentrenamiento (público) & 30 & $10^{-3}$ & Toda la red y la cabeza \\
Ajuste fino A & 5 & $10^{-3}$ & Última etapa, embedding y cabeza nueva \\
Ajuste fino B & 20 & $10^{-4}$ & Toda la red \\
Solo Polar & 40 & $10^{-3}$ & Toda la red y la cabeza \\
\bottomrule
\end{tabularx}
\end{table}

Se utiliza el optimizador AdamW con $\text{weight decay} = 10^{-4}$, planificador de coseno (\emph{CosineAnnealingLR}), tamaño de lote de 64 y semilla fija igual a 42. La detención temprana se basa en el \textbf{EER de validación} con paciencia de 5 épocas, y no en la exactitud de clasificación, porque el objetivo es separar personas que la red no ha visto. Durante el entrenamiento se aplica aumento de datos: ruido gaussiano (desviación 0.02), deformación de magnitud (0.10), deformación temporal (0.05) y deriva de línea base (amplitud 0.10).

\subsection{Partición de los datos}

Las particiones se hacen siempre \textbf{por persona} y nunca por ventana, porque las ventanas de una misma persona están correlacionadas y mezclarlas inflaría el desempeño.

\begin{table}[htbp]
\centering
\caption{Grupos de personas utilizados.}
\label{tab:particion}
\small
\begin{tabularx}{\textwidth}{l X X}
\toprule
\textbf{Grupo} & \textbf{Uso} & \textbf{Regla} \\
\midrule
Personas de ECG-ID & Preentrenamiento & Nunca se mezclan con las del Polar \\
Polar: entrenamiento & Ajuste fino o entrenamiento solo Polar & Sesiones 1 y 2 entrenan; la sesión 3 sirve como verificación entre días \\
Polar: validación & Detención temprana y umbral & Personas no vistas; nunca actualizan los pesos \\
Polar: \emph{hold-out} & Prueba final & Se usan una sola vez, al final \\
\bottomrule
\end{tabularx}
\end{table}

Con 20 voluntarios la partición prevista es de 14 personas de entrenamiento, 2 de validación y 4 de prueba (\emph{hold-out}), con un mínimo de 2 personas \emph{hold-out}. Como mejora opcional, puede usarse validación cruzada por persona en 5 pliegues de 4 personas, de modo que las 20 sirvan una vez como personas no vistas. La lista de validación y de \emph{hold-out} se elige al azar una sola vez, después de recolectar los datos y antes de entrenar con datos Polar, y se registra en el archivo de configuración; el código se detiene con un error si alguna de esas personas aparece en los datos de entrenamiento. Con 4 personas de prueba hay solo 12 combinaciones plantilla--impostor, por lo que el EER final tendrá intervalos de confianza amplios.

\subsection{Métricas de evaluación}

\begin{itemize}
  \item \textbf{Pares genuinos}: ventanas de una sesión posterior contra la plantilla de la misma persona (entre sesiones).
  \item \textbf{Pares impostores}: ventanas de una persona contra la plantilla de otra.
  \item \textbf{FAR} (tasa de falsa aceptación) y \textbf{FRR} (tasa de falso rechazo) para un umbral dado, y la \textbf{EER}, el punto en que ambas son iguales:
  \begin{equation}
    \mathrm{FAR}(\tau) = \mathrm{FRR}(\tau) = \mathrm{EER}
  \end{equation}
  \item Curva ROC y AUC.
  \item El umbral $\tau$ se elige en el punto de EER con datos de validación únicamente. Los intervalos de confianza se obtienen por \emph{bootstrap} remuestreando personas (1000 repeticiones), y cada cifra se reporta junto con el número de pares genuinos e impostores en que se basa.
\end{itemize}

\subsection{Diseño de la decisión continua}
\label{sec:decision}

La decisión final no se basa en una sola ventana, pues una lectura aislada puede fallar por un movimiento o por mal contacto de los electrodos. La regla de decisión continua es la siguiente:

\begin{enumerate}
  \item Cada 2.5~s se obtiene una ventana nueva y se calcula su similitud coseno con la plantilla del usuario.
  \item Si la similitud es mayor o igual que el umbral $\tau$, la ventana se acepta y el contador de rechazos consecutivos se reinicia a cero.
  \item Si es menor, se rechaza y el contador aumenta en uno. Mientras el contador sea menor que 3, el usuario no percibe nada.
  \item Con 3 rechazos consecutivos aparece una ventana emergente que pide al usuario sentarse derecho y permanecer quieto unos segundos, y el sistema toma una nueva muestra.
  \item Si la nueva muestra es aceptada, el examen continúa y el contador se reinicia. Si es rechazada, el examen se bloquea.
\end{enumerate}

Dado que la primera ventana completa se obtiene a los 5~s y las siguientes a intervalos de 2.5~s, el aviso al usuario aparece como mínimo unos 10~s después de que comienza una suplantación, y el bloqueo ocurre después de la nueva muestra, de otros 5~s como mínimo. Si cada rechazo injustificado ocurriera de forma independiente con probabilidad igual a la FRR, la probabilidad de 3 consecutivos sería $\mathrm{FRR}^3$; esto es solo una aproximación, pues las ventanas traslapadas no son independientes, y la tasa real de bloqueos injustificados se medirá con grabaciones completas en el Capítulo~\ref{cap:resultados}. El valor de 3 rechazos es un parámetro de diseño que puede ajustarse con los resultados. Esta lógica es independiente de la red, que solo produce los embeddings.

\subsection{Simulación y pruebas con señales sintéticas}

Antes de contar con datos reales se construyó un generador de señales ECG sintéticas para verificar el funcionamiento de cada módulo. Con ellas se comprobó que la red devuelve la forma de salida esperada $(B, 128)$ con norma igual a 1, que el resultado es determinista en modo evaluación, y que un entrenamiento pequeño sobre un lote reduce la pérdida. También se comprobó que la evaluación da un EER cercano a 0 con puntuaciones perfectamente separables y cercano a 0.5 con puntuaciones aleatorias, que un guardado y carga del modelo producen los mismos embeddings, y que el programa se detiene si una persona \emph{hold-out} aparece en el entrenamiento.

\section{Construcción e implementación del prototipo}
\label{sec:construccion}

Esta sección describe cómo el diseño se convirtió en un prototipo funcional: la organización del software, la adquisición de datos con el sensor, el protocolo de recolección, las pruebas parciales y los cambios realizados durante las iteraciones.

\subsection{Hardware}

El prototipo utiliza una banda pectoral Polar H10, que registra un canal de ECG a 130~Hz y lo transmite por Bluetooth Low Energy a una computadora. Antes de cada grabación se humedecen los electrodos de la banda para asegurar el contacto. El entrenamiento se realiza en una computadora portátil HP OmniBook X con 16~GB de RAM y CPU de 6 núcleos, sin GPU. \pendiente{agregar fotografía de la banda}

\subsection{Implementación del software}

El software se desarrolló en Python 3.14 con PyTorch 2.14 (versión para CPU), NumPy 2.3, SciPy 1.16, scikit-learn, Matplotlib 3.10 y Bleak 3.0 para la conexión Bluetooth, y se organizó de forma modular para facilitar su explicación, prueba y mantenimiento.

\begin{table}[htbp]
\centering
\caption{Módulos del software.}
\label{tab:modulos}
\small
\begin{tabularx}{\textwidth}{l X}
\toprule
\textbf{Módulo} & \textbf{Función} \\
\midrule
\texttt{config} & Todos los parámetros, la semilla y la lista de personas \emph{hold-out} \\
\texttt{record}, \texttt{check} & Conexión BLE con el Polar y chequeo rápido de la lectura antes de grabar \\
\texttt{protocol}, \texttt{session}, \texttt{status} & Protocolo de sesiones, grabación guiada con nombres automáticos de archivo y seguimiento del avance \\
\texttt{quality} & Control de calidad de la señal (señal plana, ruido, picos anómalos, número de latidos) \\
\texttt{windowing}, \texttt{data} & Filtrado, ventaneo, conjuntos de datos, particiones por persona y aumento de datos \\
\texttt{public\_data} & Carga y remuestreo de la base pública ECG-ID \\
\texttt{model} & Red neuronal de embeddings \\
\texttt{train}, \texttt{checkpoint} & Preentrenamiento, ajuste fino, entrenamiento solo Polar y guardado de modelos con archivo JSON de configuración y métricas \\
\texttt{evaluate} & EER, ROC, FAR/FRR, intervalos \emph{bootstrap} y tablas de resultados \\
\texttt{synthetic} & Señales sintéticas para pruebas \\
\bottomrule
\end{tabularx}
\end{table}

Cada etapa de entrenamiento se ejecuta con un solo comando y los modelos se guardan como versiones numeradas junto con un archivo que registra la configuración, la semilla, la lista de personas \emph{hold-out}, la versión de los datos y las métricas. Solo se almacenan los pesos del codificador.

\subsection{Protocolo de recolección de datos}

Con el Polar H10 se grabaron 20 voluntarios con identificadores anónimos fijos (P01, P02, \dots). Cada voluntario realiza \textbf{3 sesiones en días distintos}, con una separación de 2 a 3 días, para medir la variación de la señal entre días. Cada sesión dura aproximadamente 5.5 minutos y se divide en cuatro actividades:

\begin{table}[htbp]
\centering
\caption{Actividades de cada sesión de grabación.}
\label{tab:sesion}
\small
\begin{tabular}{l c}
\toprule
\textbf{Actividad} & \textbf{Duración (min)} \\
\midrule
Reposo & 2 \\
Lectura & 1.5 \\
Resolución de problemas & 1 \\
Hablando & 0.5 \\
\bottomrule
\end{tabular}
\end{table}

El programa guía al voluntario, evalúa la calidad de cada actividad y permite repetir una si salió mal. Los datos crudos no se suben al repositorio. Cada participante firmó un formato de consentimiento informado que especificaba las actividades que se realizarían, los cuidados durante la grabación, el uso que se daría a los datos (exclusivamente de investigación, con identificadores anónimos) y que la grabación no implicaba riesgos para su salud. \pendiente{elaborar y recabar el formato firmado; anexarlo al documento}

\subsection{Calibración y validación de instrumentos}

La calidad de cada actividad se verifica automáticamente con \texttt{quality}, y antes de grabar se corre un chequeo de unos 12~s con \texttt{check} que advierte si hay que ajustar la banda. Para el preentrenamiento se comprobó la frecuencia de muestreo de la base pública y se remuestreó a 130~Hz.

\subsection{Integración y pruebas parciales}

Cada módulo cuenta con pruebas automáticas con señales sintéticas (por ejemplo, forma de salida y norma del embedding, ausencia de fugas de personas \emph{hold-out} hacia el entrenamiento, evaluación con casos de EER conocido, guardado y carga de modelos, y comprobación de la calidad de señal). Primero se verificó el modelo y los datos por separado, después el entrenamiento y la evaluación, y finalmente la adquisición con el sensor.

\subsection{Iteraciones del prototipo}

\begin{table}[htbp]
\centering
\caption{Iteraciones del prototipo.}
\label{tab:iteraciones}
\small
\begin{tabularx}{\textwidth}{c X X}
\toprule
\textbf{Versión} & \textbf{Implementación principal} & \textbf{Propósito} \\
\midrule
1 & Red de embeddings, entrenamiento y evaluación; preentrenamiento con ECG-ID (modelo v1) & Comprobar que la red aprende a distinguir personas \\
2 & Adquisición con el Polar, control de calidad y sesiones guiadas de cuatro actividades & Obtener datos propios con un protocolo repetible \\
3 & Recolección con los 20 voluntarios y ajuste fino con el Polar & Adaptar la red al sensor y a las personas del estudio \\
4 & Lógica de decisión continua (3 rechazos, nueva muestra, bloqueo) y, de ser necesario, mejoras de preprocesamiento & Detectar un impostor durante el examen \\
5 & Interfaz del examen de unas 30 preguntas de comprensión lectora y pruebas de tiempo de detección & Integrar el sistema completo \\
\bottomrule
\end{tabularx}
\end{table}

\pendiente{las versiones 1 y 2 corresponden a lo ya implementado (confirmado en el historial de git); las versiones 3 a 5 son el plan de trabajo actual. Ajustar al entregar, indicando para cada versión el problema detectado, el cambio realizado y la razón técnica}

\subsection{Simulación del examen}

El examen simulado consta de unas 30 preguntas de comprensión lectora: se presentan lecturas y, después de cada una, aproximadamente 5 preguntas. Mientras el usuario lo contesta, el sistema verifica su identidad de manera continua con la lógica descrita en la Sección~\ref{sec:decision}. \pendiente{definir la plataforma, la duración y las lecturas empleadas}

\section{Presentación de la solución desarrollada}
\label{sec:solucion}

Esta sección presenta el resultado final del proceso de diseño y construcción: qué solución se obtuvo, cómo funciona y cuáles son sus características. Aquí se describe lo que se construyó; el Capítulo~\ref{cap:resultados} demuestra qué tan bien funcionó.

\subsection{Descripción de la solución}

La solución desarrollada es un sistema de autenticación continua que verifica la identidad de una persona durante un examen simulado a partir de su ECG. El usuario se registra una vez mediante una sesión de enrolamiento, de la cual se obtiene su plantilla. Durante el examen, la banda Polar H10 transmite la señal por Bluetooth; el sistema la filtra y la divide en ventanas de 5~s con 50\,\% de traslape, y una red neuronal convolucional 1D residual genera un embedding por ventana. La similitud coseno entre el embedding y la plantilla se compara con un umbral calibrado con datos de validación para aceptar o rechazar la identidad. Los rechazos aislados no se muestran al usuario; con 3 rechazos consecutivos se le pide una nueva muestra y, si esta también falla, el examen se bloquea (Sección~\ref{sec:decision}). La primera decisión está disponible a los 5~s de iniciada la adquisición y las siguientes cada 2.5~s. La interfaz del examen está por definirse.

\subsection{Secuencia de funcionamiento}

\begin{enumerate}
  \item Se coloca la banda y se verifica la calidad de la señal.
  \item El usuario realiza el registro y se calcula su plantilla.
  \item Inicia el examen y comienza la adquisición continua.
  \item Cada 2.5~s (por el traslape) se procesa una nueva ventana de 5~s.
  \item Se calcula la similitud con la plantilla y se decide si se acepta la identidad.
  \item Los rechazos aislados no se notifican. Con 3 consecutivos aparece un aviso para sentarse derecho y se toma una nueva muestra; si también es rechazada, el examen se bloquea.
\end{enumerate}

\subsection{Ficha técnica}

\begin{table}[htbp]
\centering
\caption{Ficha técnica de la solución.}
\label{tab:ficha}
\small
\begin{tabularx}{\textwidth}{l X}
\toprule
\textbf{Característica} & \textbf{Descripción} \\
\midrule
Tipo de sistema & Autenticación biométrica continua basada en ECG \\
Sensor & Polar H10, un canal, 130~Hz, conexión BLE \\
Entrada del modelo & Ventana de 650 muestras (5~s), filtrada y normalizada \\
Modelo & CNN 1D residual de 371{,}024 parámetros \\
Salida & Embedding de 128 dimensiones con norma L2 igual a 1 \\
Comparación & Similitud coseno con la plantilla del usuario \\
Entrenamiento & Preentrenamiento con ECG-ID y ajuste fino con voluntarios del Polar \\
Hardware de cómputo & CPU, sin GPU \\
Lenguaje y bibliotecas & Python, PyTorch, NumPy, SciPy, scikit-learn \\
Datos registrados & Señal por sesión con ID anónimo, calidad y metadatos \\
Aplicación & Verificación continua de identidad en un examen simulado \\
\bottomrule
\end{tabularx}
\end{table}

\subsection{Relación entre requerimientos e implementación}

\begin{table}[htbp]
\centering
\caption{Relación general entre requerimientos e implementación.}
\label{tab:req_impl}
\small
\begin{tabularx}{\textwidth}{l X}
\toprule
\textbf{Requerimiento} & \textbf{Implementación en la solución} \\
\midrule
RF-01 Adquirir ECG & Módulo \texttt{record} con el Polar H10 por BLE \\
RF-02 Calidad de la señal & Módulos \texttt{quality} y \texttt{check} \\
RF-03 Embedding por ventana & Red 1D-CNN residual del módulo \texttt{model} \\
RF-04 Plantilla del usuario & Promedio de embeddings de la sesión de enrolamiento \\
RF-05 Decisión por similitud & Similitud coseno y umbral elegido en validación \\
RD-01, RD-02 Desempeño en personas nuevas y entre días & Partición por persona y evaluación con \texttt{evaluate} \\
RF-06, RD-03, RD-04 Detección de impostor sin molestar al legítimo & Regla de 3 rechazos consecutivos, nueva muestra y bloqueo \\
RNF-01 Sin GPU & Red de 371{,}024 parámetros entrenable en CPU \\
RNF-02 Reproducibilidad & Semilla fija y configuración centralizada en \texttt{config} \\
RNF-03 Módulos separados & Red, plantillas y decisión en componentes independientes \\
RNF-04 Anonimato & IDs fijos y datos fuera del repositorio \\
\bottomrule
\end{tabularx}
\end{table}

\subsection{Limitaciones conocidas}

\begin{itemize}
  \item El número de voluntarios es reducido, por lo que el EER del conjunto \emph{hold-out} se apoya en pocas combinaciones de personas e intervalos de confianza amplios. Los resultados se consideran una prueba de concepto.
  \item La base pública ECG-ID usa derivaciones y equipo distintos al de la banda Polar, lo que genera una diferencia de dominio que el ajuste fino busca reducir.
  \item Se utiliza un solo canal de ECG y un solo modelo de sensor.
  \item Las sesiones se grabaron en un periodo corto de semanas, por lo que no se evalúan cambios de largo plazo en la señal.
  \item La actividad durante el examen simulado es principalmente sedentaria; no se evaluó movimiento intenso ni estados emocionales extremos.
  \item Las funciones de la pérdida de tripleta y ArcFace no se implementaron en esta versión.
\end{itemize}


\section*{Cierre del capítulo}

En este capítulo se definieron los requerimientos funcionales, de desempeño y no funcionales del sistema; se generaron y compararon alternativas mediante técnicas de creatividad y una matriz de decisión; y se seleccionó una arquitectura basada en una red neuronal convolucional unidimensional con bloques residuales que produce embeddings de identidad. También se describieron el diseño de la red, el preprocesamiento, la estrategia de entrenamiento y la partición de datos por persona, así como la implementación modular del prototipo y su protocolo de recolección. En el siguiente capítulo se presentan los resultados de las pruebas experimentales, el análisis de los indicadores de desempeño y la comprobación del cumplimiento de los objetivos y de la hipótesis.
 desde el documento principal.
% ============================================================

% Marca visible para lo que el equipo debe completar o confirmar.
\newcommand{\pendiente}[1]{\textbf{\textcolor{red}{[POR DEFINIR: #1]}}}

\chapter{Diseño y desarrollo de la solución}
\label{cap:diseno}

En este capítulo se describe el proceso de diseño y desarrollo de un sistema de autenticación continua basado en la señal electrocardiográfica (ECG) adquirida con una banda Polar H10. Se presentan los requerimientos de la solución, las alternativas tecnológicas consideradas, el diseño detallado del sistema, la construcción del prototipo y las características de la solución finalmente desarrollada.

El proyecto se llevó a cabo mediante una metodología de prototipado iterativo. Primero se construyó un extractor de embeddings entrenado con una base de datos pública, después se incorporó la adquisición de señal con el sensor y el protocolo de recolección con voluntarios, y finalmente se integraron los módulos de decisión continua. Cada iteración permitió identificar problemas y mejorar el sistema. Las tablas de resultados, el desempeño medido y la comprobación de la hipótesis se presentan en el Capítulo~\ref{cap:resultados}.

\section{Especificación de requerimientos}
\label{sec:requerimientos}

Esta sección establece qué debe hacer el sistema y bajo qué condiciones debe funcionar. Los requerimientos sirven como referencia para diseñar la solución y para comprobar, en el Capítulo~\ref{cap:resultados}, si el prototipo cumple con lo solicitado.

\subsection{Descripción general del sistema}

El sistema simula un examen en computadora. Mientras la persona lo contesta, una banda Polar H10 colocada en el pecho registra su ECG. Una red neuronal convierte ventanas sucesivas de la señal en vectores de características (\emph{embeddings}) y los compara con la plantilla creada en el registro inicial del usuario. De esta forma el sistema verifica de manera continua que quien contesta el examen es la misma persona que se registró, y no solo en el momento de iniciar sesión.

\subsection{Requerimientos funcionales}

\begin{itemize}
  \item Adquirir la señal ECG del Polar H10 por Bluetooth Low Energy a 130~Hz.
  \item Verificar la calidad de la señal antes y durante la grabación (señal plana, ruido excesivo, picos anómalos, latidos insuficientes).
  \item Preprocesar la señal: filtro pasa-banda de 0.5 a 40~Hz, segmentación en ventanas de 5~s con 50\,\% de traslape y normalización z-score por ventana.
  \item Generar un embedding L2-normalizado de 128 dimensiones por ventana.
  \item Crear la plantilla de un usuario durante el registro (sesión de enrolamiento) a partir del promedio de sus embeddings.
  \item Calcular la similitud coseno entre cada ventana nueva y la plantilla, y decidir si la identidad se acepta o se rechaza.
  \item Aplicar una regla de decisión continua: los rechazos aislados no se muestran al usuario; tras 3 rechazos consecutivos se le pide una nueva muestra (por ejemplo, sentarse derecho y permanecer quieto unos segundos) y, si esta tampoco es aceptada, se bloquea el examen.
  \item Entrenar y evaluar el modelo con personas que no participaron en el entrenamiento.
  \item Registrar los datos de cada sesión con identificadores anónimos, para poder repetir los análisis.
\end{itemize}

\subsection{Requerimientos no funcionales y restricciones}

\begin{itemize}
  \item El entrenamiento debe ejecutarse en una computadora sin GPU (16~GB de RAM, CPU de 6 núcleos), en un tiempo del orden de minutos a una hora.
  \item Los experimentos deben ser reproducibles: semilla fija, un único archivo de configuración y un comando por etapa de entrenamiento.
  \item La red neuronal debe permanecer separada de la base de plantillas y de la lógica de decisión, de modo que cada componente pueda sustituirse sin modificar los demás.
  \item Los datos fisiológicos deben tratarse de forma anónima y no deben publicarse en el repositorio.
  \item El modelo debe poder explicarse por completo en cada una de sus partes por los integrantes del equipo.
  \item Restricciones del proyecto: una muestra de 20 voluntarios, un solo canal de ECG, un solo sensor disponible y un plazo de desarrollo que concluye en noviembre de 2026.
\end{itemize}

\subsection{Tabla de requerimientos}

\begin{table}[htbp]
\centering
\caption{Especificación de requerimientos y criterios de aceptación.}
\label{tab:requerimientos}
\small
\begin{tabularx}{\textwidth}{l X l X}
\toprule
\textbf{Código} & \textbf{Requerimiento} & \textbf{Tipo} & \textbf{Criterio de aceptación} \\
\midrule
RF-01 & Adquirir ECG del Polar H10 a 130~Hz & Funcional & Se guardan sesiones completas sin pérdida de conexión \\
RF-02 & Evaluar la calidad de la señal & Funcional & Una sesión con señal defectuosa se detecta y se repite \\
RF-03 & Generar un embedding por ventana de 5~s & Funcional & Salida de dimensión 128 con norma L2 igual a 1 \\
RF-04 & Crear la plantilla de cada usuario & Funcional & Plantilla construida solo con la sesión de enrolamiento \\
RF-05 & Decidir identidad por similitud coseno & Funcional & Umbral fijado solo con datos de validación \\
RF-06 & Aplicar la regla de decisión continua & Funcional & Rechazos aislados ocultos; con 3 consecutivos se pide nueva muestra y, si falla, se bloquea \\
RD-01 & Distinguir personas nunca vistas por la red & Desempeño & EER medido en personas \emph{hold-out}, claramente por debajo del azar (50\,\%) y con intervalo de confianza que no lo incluya \\
RD-02 & Mantener el desempeño entre días & Desempeño & EER entre sesiones de días distintos reportado con intervalo de confianza \\
RD-03 & Detectar a un impostor durante el examen & Desempeño & Tras 3 rechazos consecutivos (unos 10~s) se solicita una nueva muestra; si esta también falla, se bloquea el examen \\
RD-04 & No interrumpir innecesariamente al usuario legítimo & Desempeño & Los rechazos aislados no se notifican; la tasa de bloqueos injustificados se reporta en el Capítulo~\ref{cap:resultados} \\
RNF-01 & Entrenar sin GPU & No funcional & Entrenamiento completo en CPU en menos de una hora \\
RNF-02 & Reproducibilidad & No funcional & Misma semilla y configuración producen el mismo modelo \\
RNF-03 & Separación entre red, plantillas y decisión & No funcional & Cada módulo se prueba de forma independiente \\
RNF-04 & Anonimato de los participantes & No funcional & Identificadores fijos tipo P01 y datos excluidos del repositorio \\
\bottomrule
\end{tabularx}
\end{table}

No se fija una meta numérica rígida de EER: con 20 participantes y un solo canal, el objetivo es demostrar la viabilidad del enfoque. El criterio es que el EER quede claramente por debajo del azar y se compare con los valores reportados en el estado del arte del Capítulo~\ref{cap:marco}. Por ello, los resultados del proyecto se consideran una prueba de concepto.

\section{Generación, evaluación y selección de alternativas}
\label{sec:alternativas}

Esta sección muestra cómo se generaron distintas propuestas para resolver el problema y por qué se eligió una de ellas. Las técnicas de creatividad no son las alternativas finales: son procedimientos para generar ideas que después se organizan, comparan y seleccionan.

% NOTA PARA EL EQUIPO: las ideas, tablas y calificaciones de esta sección son una
% PROPUESTA de partida. Deben discutirse y ajustarse según lo que el equipo realmente
% consideró, y sustentarse con literatura o pruebas preliminares. Conserven evidencias
% (fotos de la lluvia de ideas, mapa mental) para anexarlas.

\subsection{Técnicas para la generación de ideas}

\subsubsection{Lluvia de ideas}

Pregunta generadora: ¿qué sensor, qué representación de la señal y qué modelo podrían utilizarse para verificar de forma continua la identidad de una persona a partir de su ECG? Las ideas obtenidas fueron: banda Polar H10, un sensor de ECG de diseño propio, otros equipos comerciales de mayor costo (por ejemplo, Movesense o Zephyr BioHarness), señal cruda o latidos segmentados, rasgos manuales (intervalos y amplitudes de las ondas P, QRS y T), máquina de vectores de soporte (SVM), red neuronal convolucional 1D como clasificador, red 1D con embeddings, pérdida de tripleta, preentrenamiento con una base de datos pública y suavizado de las decisiones en el tiempo. \pendiente{confirmar con el equipo qué equipo comercial más caro se consideró}

\subsubsection{Mapa mental}

Se colocó en el centro el concepto <<autenticación continua con ECG>> y se crearon ramas para adquisición, preprocesamiento, representación de la señal, modelo, entrenamiento, decisión continua, interfaz y validación. El mapa permitió verificar que las propuestas contemplaran el sistema completo y no solamente la red neuronal. \pendiente{insertar figura del mapa mental}

\subsubsection{SCAMPER}

\begin{table}[htbp]
\centering
\caption{Aplicación de SCAMPER al sistema de autenticación continua.}
\label{tab:scamper}
\small
\begin{tabularx}{\textwidth}{l X X}
\toprule
\textbf{Acción} & \textbf{Pregunta aplicada} & \textbf{Ejemplo de idea} \\
\midrule
Sustituir & ¿Qué componente puede reemplazarse? & Sustituir rasgos manuales por rasgos aprendidos por una red \\
Combinar & ¿Qué elementos pueden trabajar juntos? & Combinar preentrenamiento público con ajuste fino con datos del Polar \\
Adaptar & ¿Qué solución externa puede aprovecharse? & Adaptar el aprendizaje de métricas usado en reconocimiento facial \\
Modificar & ¿Qué característica puede cambiarse? & Cambiar la longitud de la ventana de 5~s \\
Otros usos & ¿Puede un elemento cumplir otra función? & Usar la penúltima capa de un clasificador como huella de identidad \\
Eliminar & ¿Qué elemento no es indispensable? & Eliminar la detección explícita de picos R y trabajar sobre la señal filtrada \\
Reordenar & ¿Qué ocurre al cambiar el orden? & Suavizar decisiones en el tiempo en lugar de decidir por ventana aislada \\
\bottomrule
\end{tabularx}
\end{table}

\subsubsection{Combinación de elementos y matriz morfológica}

La matriz morfológica combina opciones de distintos subsistemas para construir alternativas completas y técnicamente coherentes.

\begin{table}[htbp]
\centering
\caption{Matriz morfológica del sistema.}
\label{tab:morfologica}
\small
\begin{tabularx}{\textwidth}{l X X X}
\toprule
\textbf{Función} & \textbf{Opción 1} & \textbf{Opción 2} & \textbf{Opción 3} \\
\midrule
Adquisición & Polar H10 & Sensor de ECG de diseño propio & Equipo comercial de mayor costo \\
Entrada del modelo & Rasgos manuales & Ventana cruda filtrada & Latidos segmentados \\
Modelo & SVM & CNN 1D como clasificador & CNN 1D con embeddings \\
Pérdida & Entropía cruzada & Tripleta & ArcFace \\
Datos de entrenamiento & Solo Polar & Preentrenado y ajustado & Solo base pública \\
Decisión & Por ventana & Promedio móvil & Acumulación de evidencia \\
\bottomrule
\end{tabularx}
\end{table}

\subsection{Alternativas de solución obtenidas}

\begin{table}[htbp]
\centering
\caption{Alternativas conceptuales.}
\label{tab:alternativas}
\small
\begin{tabularx}{\textwidth}{l X X X}
\toprule
\textbf{Alternativa} & \textbf{Configuración} & \textbf{Ventaja principal} & \textbf{Limitación principal} \\
\midrule
A: clásica & Polar H10, rasgos manuales, SVM, decisión por ventana & Pocos datos necesarios y fácil de explicar & Depende de detectar con precisión las ondas; menor capacidad de generalizar a personas nuevas \\
B: embeddings con preentrenamiento & Polar H10, ventana filtrada, CNN 1D residual con embeddings, preentrenamiento público y ajuste fino, entropía cruzada & Aprende rasgos automáticamente; permite registrar personas nuevas sin reentrenar; el preentrenamiento compensa pocos voluntarios & Diferencia entre la base pública y el Polar; requiere más tiempo de entrenamiento \\
C: aprendizaje de métricas avanzado & Igual que B pero con pérdida de tripleta o ArcFace y detección de picos R & Mayor separación potencial entre identidades & Mayor complejidad y sensibilidad a hiperparámetros con pocos datos \\
\bottomrule
\end{tabularx}
\end{table}

\subsection{Criterios de evaluación}

\begin{table}[htbp]
\centering
\caption{Criterios de evaluación y ponderación.}
\label{tab:criterios}
\small
\begin{tabularx}{\textwidth}{l l X}
\toprule
\textbf{Criterio} & \textbf{Ponderación} & \textbf{Justificación} \\
\midrule
Capacidad de generalizar a personas nuevas & 30\,\% & El sistema debe funcionar con usuarios que la red nunca vio \\
Viabilidad con pocos datos & 20\,\% & Solo se cuenta con 20 voluntarios, una cantidad pequeña para entrenar redes profundas \\
Costo computacional & 15\,\% & Solo se dispone de CPU \\
Facilidad de implementación y explicación & 20\,\% & El equipo debe dominar cada componente y cumplir el plazo \\
Aprovechamiento del sensor disponible & 15\,\% & El Polar H10 ya está en poder del equipo \\
\bottomrule
\end{tabularx}
\end{table}

Cada alternativa se califica del 1 al 5, donde 1 es un desempeño muy deficiente y 5 uno excelente. El resultado ponderado se obtiene multiplicando la calificación por la ponderación; la suma permite comparar las alternativas.

\begin{table}[htbp]
\centering
\caption{Matriz de decisión (calificaciones preliminares, por justificar).}
\label{tab:decision}
\small
\begin{tabular}{l c c c c}
\toprule
\textbf{Criterio} & \textbf{Pond.} & \textbf{A} & \textbf{B} & \textbf{C} \\
\midrule
Generalización a personas nuevas & 0.30 & 2 & 4 & 5 \\
Viabilidad con pocos datos & 0.20 & 4 & 4 & 2 \\
Costo computacional & 0.15 & 5 & 4 & 3 \\
Facilidad de implementación y explicación & 0.20 & 4 & 4 & 2 \\
Aprovechamiento del sensor & 0.15 & 4 & 4 & 4 \\
\midrule
\textbf{Total ponderado} & 1.00 & 3.55 & 4.00 & 3.35 \\
\bottomrule
\end{tabular}
\end{table}

\pendiente{justificar cada calificación con literatura, pruebas preliminares o cálculos; las cifras mostradas son un borrador}

\subsection{Selección y justificación}

Después de aplicar las técnicas de creatividad se obtuvieron tres alternativas de solución, evaluadas según su capacidad de generalizar a personas nuevas, su viabilidad con pocos datos, su costo computacional, su facilidad de implementación y el aprovechamiento del sensor disponible. La alternativa B obtuvo la mayor puntuación porque ofrece un equilibrio adecuado entre capacidad de generalización y viabilidad: el uso de embeddings permite registrar usuarios nuevos solo calculando su plantilla, y el preentrenamiento con una base de datos pública mitiga la escasez de voluntarios. La alternativa C se consideró una mejora futura en caso de que el error obtenido con B no fuera suficiente, por lo que la pérdida de tripleta y ArcFace se dejaron deliberadamente fuera de esta versión. En cuanto al sensor, se consideró construir uno propio, pero se encontró que el Polar H10 es más preciso, por lo que se descartó. Los equipos comerciales de mayor costo no ofrecían una ventaja que justificara su precio, y el Polar H10 cuenta con estudios que lo validan frente a un ECG de referencia. Por lo anterior, se seleccionó la alternativa B, con el Polar H10, como base para el diseño detallado.

\section{Modelado, cálculos, simulación y diseño detallado}
\label{sec:diseno}

Esta sección explica técnicamente cómo se diseñó la alternativa seleccionada antes de construirla, de modo que otra persona pueda comprender la arquitectura del sistema y reproducir el diseño.

\subsection{Arquitectura general}

El sistema se organiza como una cadena de bloques independientes:
\begin{center}
\emph{Polar H10} $\rightarrow$ adquisición BLE $\rightarrow$ control de calidad $\rightarrow$ filtrado y ventaneo $\rightarrow$ red de embeddings $\rightarrow$ similitud coseno con la plantilla $\rightarrow$ decisión
\end{center}
\pendiente{insertar diagrama de bloques}

La red (los pesos) se mantiene separada de la base de plantillas y de la lógica de decisión. Así, registrar a un usuario nuevo solo requiere calcular su plantilla, sin reentrenar la red.

\subsection{Adquisición y preprocesamiento de la señal}

La banda Polar H10 entrega el ECG a $f_s = 130$~Hz. La señal se procesa de la siguiente forma:
\begin{enumerate}
  \item \textbf{Filtro pasa-banda} de 0.5 a 40~Hz, que elimina la deriva de la línea base y atenúa el ruido muscular de alta frecuencia. La frecuencia de Nyquist es de 65~Hz, por lo que la banda es válida.
  \item \textbf{Ventaneo} en segmentos de 5~s, es decir, $5 \times 130 = 650$ muestras, con 50\,\% de traslape.
  \item \textbf{Normalización z-score por ventana}, para eliminar diferencias de amplitud debidas al contacto de los electrodos:
  \begin{equation}
    \tilde{x}_i = \frac{x_i - \mu}{\sigma}
  \end{equation}
  donde $\mu$ y $\sigma$ son la media y la desviación estándar de la ventana.
\end{enumerate}
Cada ventana es un tensor de forma $(1, 650)$. El sistema no detecta explícitamente los picos R para el modelo: la red extrae por sí misma las características morfológicas y temporales.

\subsection{Red neuronal de embeddings}

\subsubsection{Tipo de red}

Se utiliza una red neuronal convolucional unidimensional (1D-CNN) con bloques residuales, usada como \textbf{extractor de embeddings} (aprendizaje de métricas) y no como clasificador final. Su salida es un vector de 128 dimensiones con norma L2 igual a 1, de modo que la similitud coseno entre dos ventanas indica si provienen de la misma persona.

\subsubsection{Arquitectura capa por capa}

\begin{table}[htbp]
\centering
\caption{Arquitectura de la red para una ventana de entrada $(1, 650)$.}
\label{tab:arquitectura}
\small
\begin{tabularx}{\textwidth}{c l X c c c}
\toprule
\textbf{\#} & \textbf{Bloque} & \textbf{Operación} & \textbf{Kernel / stride} & \textbf{Canales} & \textbf{Salida} \\
\midrule
0 & Entrada & Ventana z-score & -- & 1 & $1 \times 650$ \\
1 & Stem & Conv1d + BatchNorm + ReLU & 7 / 1 & $1 \rightarrow 16$ & $16 \times 650$ \\
2 & Etapa 1 & Bloque residual & 5 / 2 & $16 \rightarrow 32$ & $32 \times 325$ \\
3 & Etapa 2 & Bloque residual & 5 / 2 & $32 \rightarrow 64$ & $64 \times 163$ \\
4 & Etapa 3 & Bloque residual & 5 / 2 & $64 \rightarrow 128$ & $128 \times 82$ \\
5 & Etapa 4 & Bloque residual & 5 / 2 & $128 \rightarrow 128$ & $128 \times 41$ \\
6 & Pooling global & Promedio sobre el tiempo & -- & 128 & 128 \\
7 & Embedding & Capa densa & -- & $128 \rightarrow 128$ & 128 \\
8 & Normalización & L2 & -- & 128 & 128 (norma 1) \\
9 & Cabeza & Capa densa sobre el embedding sin normalizar (solo entrenamiento) & -- & $128 \rightarrow N$ & $N$ \\
\bottomrule
\end{tabularx}
\end{table}

Todas las convoluciones usan \texttt{bias=False}, pues la normalización por lotes (BatchNorm) ya aporta el desplazamiento. La red tiene \textbf{371{,}024 parámetros} en el codificador; la cabeza agrega $128N + N$ parámetros, con $N$ igual al número de personas de entrenamiento, y se descarta al terminar.

\subsubsection{Bloque residual}

Cada etapa usa el mismo bloque. La rama principal aplica una convolución con $k=5$ y \emph{stride} 2, BatchNorm y ReLU, y después una segunda convolución con $k=5$ y \emph{stride} 1 seguida de BatchNorm. El atajo (\emph{skip connection}) es una convolución de $1\times1$ con \emph{stride} 2 seguida de BatchNorm, para que las dimensiones coincidan en la suma. La ReLU se aplica después de sumar ambas ramas:
\begin{equation}
  \mathbf{y} = \mathrm{ReLU}\big(F(\mathbf{x}) + S(\mathbf{x})\big)
\end{equation}
donde $F$ es la rama principal y $S$ el atajo. El \emph{stride} 2 reduce a la mitad la longitud temporal en cada etapa ($650 \rightarrow 325 \rightarrow 163 \rightarrow 82 \rightarrow 41$); con ello el campo receptivo de la última etapa abarca varios latidos.

\subsubsection{Embedding y similitud}

El pooling global promedio hace que la huella sea independiente de en qué punto de la ventana cae cada latido. La salida se normaliza:
\begin{equation}
  \mathbf{e} = \frac{\mathbf{z}}{\lVert \mathbf{z} \rVert_2}
\end{equation}
y la similitud entre una ventana con embedding $\mathbf{e}$ y una plantilla $\mathbf{t}$ es el coseno del ángulo entre ambos, que al estar normalizados equivale al producto punto:
\begin{equation}
  s(\mathbf{e}, \mathbf{t}) = \mathbf{e} \cdot \mathbf{t}
\end{equation}
La plantilla de una persona es el promedio de los embeddings de su sesión de enrolamiento. La identidad se acepta si $s$ supera un umbral $\tau$.

\subsection{Estrategia de entrenamiento}

\begin{enumerate}
  \item \textbf{Preentrenamiento} con la base pública ECG-ID, remuestreada a 130~Hz: la red se entrena como clasificador de identidades con entropía cruzada y \emph{label smoothing} de 0.1. Su propósito es aprender rasgos generales del ECG.
  \item \textbf{Ajuste fino} con los voluntarios del Polar: se reemplaza la cabeza por una nueva. En la etapa A se congelan el stem y las tres primeras etapas (con su BatchNorm en modo evaluación) y se entrenan la última etapa, el embedding y la cabeza; en la etapa B se descongela toda la red con una tasa de aprendizaje diez veces menor.
  \item \textbf{Alternativa solo Polar}: mismo procedimiento sin preentrenamiento, seleccionado con una bandera de configuración. Permite medir si el preentrenamiento realmente ayuda.
\end{enumerate}

\begin{table}[htbp]
\centering
\caption{Hiperparámetros por etapa de entrenamiento.}
\label{tab:hiperparametros}
\small
\begin{tabularx}{\textwidth}{X c c X}
\toprule
\textbf{Etapa} & \textbf{Épocas} & \textbf{Tasa de aprendizaje} & \textbf{Qué se entrena} \\
\midrule
Preentrenamiento (público) & 30 & $10^{-3}$ & Toda la red y la cabeza \\
Ajuste fino A & 5 & $10^{-3}$ & Última etapa, embedding y cabeza nueva \\
Ajuste fino B & 20 & $10^{-4}$ & Toda la red \\
Solo Polar & 40 & $10^{-3}$ & Toda la red y la cabeza \\
\bottomrule
\end{tabularx}
\end{table}

Se utiliza el optimizador AdamW con $\text{weight decay} = 10^{-4}$, planificador de coseno (\emph{CosineAnnealingLR}), tamaño de lote de 64 y semilla fija igual a 42. La detención temprana se basa en el \textbf{EER de validación} con paciencia de 5 épocas, y no en la exactitud de clasificación, porque el objetivo es separar personas que la red no ha visto. Durante el entrenamiento se aplica aumento de datos: ruido gaussiano (desviación 0.02), deformación de magnitud (0.10), deformación temporal (0.05) y deriva de línea base (amplitud 0.10).

\subsection{Partición de los datos}

Las particiones se hacen siempre \textbf{por persona} y nunca por ventana, porque las ventanas de una misma persona están correlacionadas y mezclarlas inflaría el desempeño.

\begin{table}[htbp]
\centering
\caption{Grupos de personas utilizados.}
\label{tab:particion}
\small
\begin{tabularx}{\textwidth}{l X X}
\toprule
\textbf{Grupo} & \textbf{Uso} & \textbf{Regla} \\
\midrule
Personas de ECG-ID & Preentrenamiento & Nunca se mezclan con las del Polar \\
Polar: entrenamiento & Ajuste fino o entrenamiento solo Polar & Sesiones 1 y 2 entrenan; la sesión 3 sirve como verificación entre días \\
Polar: validación & Detención temprana y umbral & Personas no vistas; nunca actualizan los pesos \\
Polar: \emph{hold-out} & Prueba final & Se usan una sola vez, al final \\
\bottomrule
\end{tabularx}
\end{table}

Con 20 voluntarios la partición prevista es de 14 personas de entrenamiento, 2 de validación y 4 de prueba (\emph{hold-out}), con un mínimo de 2 personas \emph{hold-out}. Como mejora opcional, puede usarse validación cruzada por persona en 5 pliegues de 4 personas, de modo que las 20 sirvan una vez como personas no vistas. La lista de validación y de \emph{hold-out} se elige al azar una sola vez, después de recolectar los datos y antes de entrenar con datos Polar, y se registra en el archivo de configuración; el código se detiene con un error si alguna de esas personas aparece en los datos de entrenamiento. Con 4 personas de prueba hay solo 12 combinaciones plantilla--impostor, por lo que el EER final tendrá intervalos de confianza amplios.

\subsection{Métricas de evaluación}

\begin{itemize}
  \item \textbf{Pares genuinos}: ventanas de una sesión posterior contra la plantilla de la misma persona (entre sesiones).
  \item \textbf{Pares impostores}: ventanas de una persona contra la plantilla de otra.
  \item \textbf{FAR} (tasa de falsa aceptación) y \textbf{FRR} (tasa de falso rechazo) para un umbral dado, y la \textbf{EER}, el punto en que ambas son iguales:
  \begin{equation}
    \mathrm{FAR}(\tau) = \mathrm{FRR}(\tau) = \mathrm{EER}
  \end{equation}
  \item Curva ROC y AUC.
  \item El umbral $\tau$ se elige en el punto de EER con datos de validación únicamente. Los intervalos de confianza se obtienen por \emph{bootstrap} remuestreando personas (1000 repeticiones), y cada cifra se reporta junto con el número de pares genuinos e impostores en que se basa.
\end{itemize}

\subsection{Diseño de la decisión continua}
\label{sec:decision}

La decisión final no se basa en una sola ventana, pues una lectura aislada puede fallar por un movimiento o por mal contacto de los electrodos. La regla de decisión continua es la siguiente:

\begin{enumerate}
  \item Cada 2.5~s se obtiene una ventana nueva y se calcula su similitud coseno con la plantilla del usuario.
  \item Si la similitud es mayor o igual que el umbral $\tau$, la ventana se acepta y el contador de rechazos consecutivos se reinicia a cero.
  \item Si es menor, se rechaza y el contador aumenta en uno. Mientras el contador sea menor que 3, el usuario no percibe nada.
  \item Con 3 rechazos consecutivos aparece una ventana emergente que pide al usuario sentarse derecho y permanecer quieto unos segundos, y el sistema toma una nueva muestra.
  \item Si la nueva muestra es aceptada, el examen continúa y el contador se reinicia. Si es rechazada, el examen se bloquea.
\end{enumerate}

Dado que la primera ventana completa se obtiene a los 5~s y las siguientes a intervalos de 2.5~s, el aviso al usuario aparece como mínimo unos 10~s después de que comienza una suplantación, y el bloqueo ocurre después de la nueva muestra, de otros 5~s como mínimo. Si cada rechazo injustificado ocurriera de forma independiente con probabilidad igual a la FRR, la probabilidad de 3 consecutivos sería $\mathrm{FRR}^3$; esto es solo una aproximación, pues las ventanas traslapadas no son independientes, y la tasa real de bloqueos injustificados se medirá con grabaciones completas en el Capítulo~\ref{cap:resultados}. El valor de 3 rechazos es un parámetro de diseño que puede ajustarse con los resultados. Esta lógica es independiente de la red, que solo produce los embeddings.

\subsection{Simulación y pruebas con señales sintéticas}

Antes de contar con datos reales se construyó un generador de señales ECG sintéticas para verificar el funcionamiento de cada módulo. Con ellas se comprobó que la red devuelve la forma de salida esperada $(B, 128)$ con norma igual a 1, que el resultado es determinista en modo evaluación, y que un entrenamiento pequeño sobre un lote reduce la pérdida. También se comprobó que la evaluación da un EER cercano a 0 con puntuaciones perfectamente separables y cercano a 0.5 con puntuaciones aleatorias, que un guardado y carga del modelo producen los mismos embeddings, y que el programa se detiene si una persona \emph{hold-out} aparece en el entrenamiento.

\section{Construcción e implementación del prototipo}
\label{sec:construccion}

Esta sección describe cómo el diseño se convirtió en un prototipo funcional: la organización del software, la adquisición de datos con el sensor, el protocolo de recolección, las pruebas parciales y los cambios realizados durante las iteraciones.

\subsection{Hardware}

El prototipo utiliza una banda pectoral Polar H10, que registra un canal de ECG a 130~Hz y lo transmite por Bluetooth Low Energy a una computadora. Antes de cada grabación se humedecen los electrodos de la banda para asegurar el contacto. El entrenamiento se realiza en una computadora portátil HP OmniBook X con 16~GB de RAM y CPU de 6 núcleos, sin GPU. \pendiente{agregar fotografía de la banda}

\subsection{Implementación del software}

El software se desarrolló en Python 3.14 con PyTorch 2.14 (versión para CPU), NumPy 2.3, SciPy 1.16, scikit-learn, Matplotlib 3.10 y Bleak 3.0 para la conexión Bluetooth, y se organizó de forma modular para facilitar su explicación, prueba y mantenimiento.

\begin{table}[htbp]
\centering
\caption{Módulos del software.}
\label{tab:modulos}
\small
\begin{tabularx}{\textwidth}{l X}
\toprule
\textbf{Módulo} & \textbf{Función} \\
\midrule
\texttt{config} & Todos los parámetros, la semilla y la lista de personas \emph{hold-out} \\
\texttt{record}, \texttt{check} & Conexión BLE con el Polar y chequeo rápido de la lectura antes de grabar \\
\texttt{protocol}, \texttt{session}, \texttt{status} & Protocolo de sesiones, grabación guiada con nombres automáticos de archivo y seguimiento del avance \\
\texttt{quality} & Control de calidad de la señal (señal plana, ruido, picos anómalos, número de latidos) \\
\texttt{windowing}, \texttt{data} & Filtrado, ventaneo, conjuntos de datos, particiones por persona y aumento de datos \\
\texttt{public\_data} & Carga y remuestreo de la base pública ECG-ID \\
\texttt{model} & Red neuronal de embeddings \\
\texttt{train}, \texttt{checkpoint} & Preentrenamiento, ajuste fino, entrenamiento solo Polar y guardado de modelos con archivo JSON de configuración y métricas \\
\texttt{evaluate} & EER, ROC, FAR/FRR, intervalos \emph{bootstrap} y tablas de resultados \\
\texttt{synthetic} & Señales sintéticas para pruebas \\
\bottomrule
\end{tabularx}
\end{table}

Cada etapa de entrenamiento se ejecuta con un solo comando y los modelos se guardan como versiones numeradas junto con un archivo que registra la configuración, la semilla, la lista de personas \emph{hold-out}, la versión de los datos y las métricas. Solo se almacenan los pesos del codificador.

\subsection{Protocolo de recolección de datos}

Con el Polar H10 se grabaron 20 voluntarios con identificadores anónimos fijos (P01, P02, \dots). Cada voluntario realiza \textbf{3 sesiones en días distintos}, con una separación de 2 a 3 días, para medir la variación de la señal entre días. Cada sesión dura aproximadamente 5.5 minutos y se divide en cuatro actividades:

\begin{table}[htbp]
\centering
\caption{Actividades de cada sesión de grabación.}
\label{tab:sesion}
\small
\begin{tabular}{l c}
\toprule
\textbf{Actividad} & \textbf{Duración (min)} \\
\midrule
Reposo & 2 \\
Lectura & 1.5 \\
Resolución de problemas & 1 \\
Hablando & 0.5 \\
\bottomrule
\end{tabular}
\end{table}

El programa guía al voluntario, evalúa la calidad de cada actividad y permite repetir una si salió mal. Los datos crudos no se suben al repositorio. Cada participante firmó un formato de consentimiento informado que especificaba las actividades que se realizarían, los cuidados durante la grabación, el uso que se daría a los datos (exclusivamente de investigación, con identificadores anónimos) y que la grabación no implicaba riesgos para su salud. \pendiente{elaborar y recabar el formato firmado; anexarlo al documento}

\subsection{Calibración y validación de instrumentos}

La calidad de cada actividad se verifica automáticamente con \texttt{quality}, y antes de grabar se corre un chequeo de unos 12~s con \texttt{check} que advierte si hay que ajustar la banda. Para el preentrenamiento se comprobó la frecuencia de muestreo de la base pública y se remuestreó a 130~Hz.

\subsection{Integración y pruebas parciales}

Cada módulo cuenta con pruebas automáticas con señales sintéticas (por ejemplo, forma de salida y norma del embedding, ausencia de fugas de personas \emph{hold-out} hacia el entrenamiento, evaluación con casos de EER conocido, guardado y carga de modelos, y comprobación de la calidad de señal). Primero se verificó el modelo y los datos por separado, después el entrenamiento y la evaluación, y finalmente la adquisición con el sensor.

\subsection{Iteraciones del prototipo}

\begin{table}[htbp]
\centering
\caption{Iteraciones del prototipo.}
\label{tab:iteraciones}
\small
\begin{tabularx}{\textwidth}{c X X}
\toprule
\textbf{Versión} & \textbf{Implementación principal} & \textbf{Propósito} \\
\midrule
1 & Red de embeddings, entrenamiento y evaluación; preentrenamiento con ECG-ID (modelo v1) & Comprobar que la red aprende a distinguir personas \\
2 & Adquisición con el Polar, control de calidad y sesiones guiadas de cuatro actividades & Obtener datos propios con un protocolo repetible \\
3 & Recolección con los 20 voluntarios y ajuste fino con el Polar & Adaptar la red al sensor y a las personas del estudio \\
4 & Lógica de decisión continua (3 rechazos, nueva muestra, bloqueo) y, de ser necesario, mejoras de preprocesamiento & Detectar un impostor durante el examen \\
5 & Interfaz del examen de unas 30 preguntas de comprensión lectora y pruebas de tiempo de detección & Integrar el sistema completo \\
\bottomrule
\end{tabularx}
\end{table}

\pendiente{las versiones 1 y 2 corresponden a lo ya implementado (confirmado en el historial de git); las versiones 3 a 5 son el plan de trabajo actual. Ajustar al entregar, indicando para cada versión el problema detectado, el cambio realizado y la razón técnica}

\subsection{Simulación del examen}

El examen simulado consta de unas 30 preguntas de comprensión lectora: se presentan lecturas y, después de cada una, aproximadamente 5 preguntas. Mientras el usuario lo contesta, el sistema verifica su identidad de manera continua con la lógica descrita en la Sección~\ref{sec:decision}. \pendiente{definir la plataforma, la duración y las lecturas empleadas}

\section{Presentación de la solución desarrollada}
\label{sec:solucion}

Esta sección presenta el resultado final del proceso de diseño y construcción: qué solución se obtuvo, cómo funciona y cuáles son sus características. Aquí se describe lo que se construyó; el Capítulo~\ref{cap:resultados} demuestra qué tan bien funcionó.

\subsection{Descripción de la solución}

La solución desarrollada es un sistema de autenticación continua que verifica la identidad de una persona durante un examen simulado a partir de su ECG. El usuario se registra una vez mediante una sesión de enrolamiento, de la cual se obtiene su plantilla. Durante el examen, la banda Polar H10 transmite la señal por Bluetooth; el sistema la filtra y la divide en ventanas de 5~s con 50\,\% de traslape, y una red neuronal convolucional 1D residual genera un embedding por ventana. La similitud coseno entre el embedding y la plantilla se compara con un umbral calibrado con datos de validación para aceptar o rechazar la identidad. Los rechazos aislados no se muestran al usuario; con 3 rechazos consecutivos se le pide una nueva muestra y, si esta también falla, el examen se bloquea (Sección~\ref{sec:decision}). La primera decisión está disponible a los 5~s de iniciada la adquisición y las siguientes cada 2.5~s. La interfaz del examen está por definirse.

\subsection{Secuencia de funcionamiento}

\begin{enumerate}
  \item Se coloca la banda y se verifica la calidad de la señal.
  \item El usuario realiza el registro y se calcula su plantilla.
  \item Inicia el examen y comienza la adquisición continua.
  \item Cada 2.5~s (por el traslape) se procesa una nueva ventana de 5~s.
  \item Se calcula la similitud con la plantilla y se decide si se acepta la identidad.
  \item Los rechazos aislados no se notifican. Con 3 consecutivos aparece un aviso para sentarse derecho y se toma una nueva muestra; si también es rechazada, el examen se bloquea.
\end{enumerate}

\subsection{Ficha técnica}

\begin{table}[htbp]
\centering
\caption{Ficha técnica de la solución.}
\label{tab:ficha}
\small
\begin{tabularx}{\textwidth}{l X}
\toprule
\textbf{Característica} & \textbf{Descripción} \\
\midrule
Tipo de sistema & Autenticación biométrica continua basada en ECG \\
Sensor & Polar H10, un canal, 130~Hz, conexión BLE \\
Entrada del modelo & Ventana de 650 muestras (5~s), filtrada y normalizada \\
Modelo & CNN 1D residual de 371{,}024 parámetros \\
Salida & Embedding de 128 dimensiones con norma L2 igual a 1 \\
Comparación & Similitud coseno con la plantilla del usuario \\
Entrenamiento & Preentrenamiento con ECG-ID y ajuste fino con voluntarios del Polar \\
Hardware de cómputo & CPU, sin GPU \\
Lenguaje y bibliotecas & Python, PyTorch, NumPy, SciPy, scikit-learn \\
Datos registrados & Señal por sesión con ID anónimo, calidad y metadatos \\
Aplicación & Verificación continua de identidad en un examen simulado \\
\bottomrule
\end{tabularx}
\end{table}

\subsection{Relación entre requerimientos e implementación}

\begin{table}[htbp]
\centering
\caption{Relación general entre requerimientos e implementación.}
\label{tab:req_impl}
\small
\begin{tabularx}{\textwidth}{l X}
\toprule
\textbf{Requerimiento} & \textbf{Implementación en la solución} \\
\midrule
RF-01 Adquirir ECG & Módulo \texttt{record} con el Polar H10 por BLE \\
RF-02 Calidad de la señal & Módulos \texttt{quality} y \texttt{check} \\
RF-03 Embedding por ventana & Red 1D-CNN residual del módulo \texttt{model} \\
RF-04 Plantilla del usuario & Promedio de embeddings de la sesión de enrolamiento \\
RF-05 Decisión por similitud & Similitud coseno y umbral elegido en validación \\
RD-01, RD-02 Desempeño en personas nuevas y entre días & Partición por persona y evaluación con \texttt{evaluate} \\
RF-06, RD-03, RD-04 Detección de impostor sin molestar al legítimo & Regla de 3 rechazos consecutivos, nueva muestra y bloqueo \\
RNF-01 Sin GPU & Red de 371{,}024 parámetros entrenable en CPU \\
RNF-02 Reproducibilidad & Semilla fija y configuración centralizada en \texttt{config} \\
RNF-03 Módulos separados & Red, plantillas y decisión en componentes independientes \\
RNF-04 Anonimato & IDs fijos y datos fuera del repositorio \\
\bottomrule
\end{tabularx}
\end{table}

\subsection{Limitaciones conocidas}

\begin{itemize}
  \item El número de voluntarios es reducido, por lo que el EER del conjunto \emph{hold-out} se apoya en pocas combinaciones de personas e intervalos de confianza amplios. Los resultados se consideran una prueba de concepto.
  \item La base pública ECG-ID usa derivaciones y equipo distintos al de la banda Polar, lo que genera una diferencia de dominio que el ajuste fino busca reducir.
  \item Se utiliza un solo canal de ECG y un solo modelo de sensor.
  \item Las sesiones se grabaron en un periodo corto de semanas, por lo que no se evalúan cambios de largo plazo en la señal.
  \item La actividad durante el examen simulado es principalmente sedentaria; no se evaluó movimiento intenso ni estados emocionales extremos.
  \item Las funciones de la pérdida de tripleta y ArcFace no se implementaron en esta versión.
\end{itemize}


\section*{Cierre del capítulo}

En este capítulo se definieron los requerimientos funcionales, de desempeño y no funcionales del sistema; se generaron y compararon alternativas mediante técnicas de creatividad y una matriz de decisión; y se seleccionó una arquitectura basada en una red neuronal convolucional unidimensional con bloques residuales que produce embeddings de identidad. También se describieron el diseño de la red, el preprocesamiento, la estrategia de entrenamiento y la partición de datos por persona, así como la implementación modular del prototipo y su protocolo de recolección. En el siguiente capítulo se presentan los resultados de las pruebas experimentales, el análisis de los indicadores de desempeño y la comprobación del cumplimiento de los objetivos y de la hipótesis.
